% Transcription — fonds Alexandre Grothendieck, shelfmark 1, batch 3 (pages 41-60).
% Established from the facsimile archives/batches/1/batch-03.pdf (a mirror of
% https://grothendieck.umontpellier.fr/1.pdf), per the skill
% transcribe-grothendieck.
%
% Pass: Fable 5.1 (claude-fable-5-1), 2026-09-06 — first pass, unchecked against
% the pages by a human.
%
% Eighteen of the twenty pages are transcribed. Pages 42 and 46 are blank and
% absent.
%
% \page{} numbers the position in the folder, as batches 1 and 2 did: the
% archivists' pencil numbers run two ahead, so page 41 here is pencilled 43
% and page 60 is pencilled 62. The note at page 3 of batch 1 records this
% once for the folder.
%
% The batch falls into three runs, all in the same fast blue-black hand on
% yellowed sheets, and the connective prose is largely reduced to dashes —
% the apparatus says so rather than inventing. The statements and the
% numbered formulas come through; the sentences between them are
% \uncertain{} word by word or \ill{}.
%
%   Page 41 — the end of the « II » of the set that closed batch 2 (pages
%     37-40): the dual of L_0(E,F) is L^1(F,E) under (u,v) = Tr vu, proved
%     through the polar decomposition v = Ww of the operator a continuous
%     form defines, then a sketch of « une méthode plus élémentaire » for
%     L^p and L^q.
%   Pages 43-57 — a run he heads « 2. Produits tensoriels hilbertiens »,
%     after a plan line « 1. Rappels et notations. Opérateurs de Fredholm »:
%     the conjugate space Ē, (1) Ē̄ = E, sesquilinear maps as linear maps on
%     E ⊗ F̄, (2) \overline{E⊗F} = Ē⊗F̄, (3) E' = Ē; the form (5)
%     (x₁⊗y₁, x₂⊗y₂) = (x₁,x₂)(y₁,y₂), (6) its orthonormality, the
%     definition of the Hilbert tensor product E ⊗^{(2)} F, Prop 1 (e_i⊗f_j
%     is an orthonormal basis), (7) ‖x⊗y‖₂ = ‖x‖‖y‖, (8); the example
%     L²_E(μ) with (9)-(10), Théorème 1 (11) L²(μ) ⊗^{(2)} E = L²_E(μ) and its
%     corollaries (12) l²(I) ⊗^{(2)} E = l²_E(I), (13) L²(μ) ⊗^{(2)} L²(ν) =
%     L²(μ⊗ν); then (14) (ā⊗b)·x = (x,a)·b, (15) (ux,y) = (u, x̄⊗y), the
%     embedding Ē ⊗^{(2)} F → L(E,F), Définition 2 of the Hilbert-Schmidt
%     operators L^{(2)}(E,F), (16) ‖u‖_∞ ≤ ‖u‖₂, Théorème 2 with (17)
%     ‖Au‖₂ ≤ ‖A‖‖u‖₂ and (18), Th. 3 with (19) ‖u‖₂ = ‖u*‖₂ = (Σρ_i²)^{1/2}
%     and its proof through u = Uh; a second « Th 3 » on Fredholm operators
%     with (20) ‖u‖₁ = ‖h‖₁ = Σρ_i, Cor. 1 (21) u = Σρ_i ā_i⊗b_i, (22) ‖u‖₁ =
%     Tr √(u*u) = Tr √(uu*), Cor. 2 (23) ‖u‖₁ = sup |Tr Au|; Th. 4 (24)
%     ‖vu‖₁ ≤ ‖v‖₂‖u‖₂ and (25) (u,v) = Tr uv* = Tr v*u with its proof,
%     Cor. 1 (26) ‖u‖₂ = sup |Tr vu|, Cor. 2 (Fredholm = product of two
%     Hilbert-Schmidt).
%   Pages 58-60 — « 3. Produits extérieurs hilbertiens. Applications : des
%     inégalités diverses », his heading: the n-fold Hilbert tensor product
%     with (1) and (2) ‖u₁⊗…⊗u_n‖ ≤ ‖u₁‖…‖u_n‖, reduced to (17) through the
%     associativity noted in the margin; the antisymmetriser (3) a_n =
%     (1/n!) Σ ε_σ σ, the definition of Λ^{n (2)} E as its image, and the
%     formula (a₁∧…∧a_n, b₁∧…∧b_n) = n!(a_n A, a_n B) = det((a_i,b_j)).
%     The numbering of formulas restarts at (1) with this section.
%
% Cancellations: pages 41, 43, 44, 47, 51, 52, 53, 54, 59 and 60 carry
% lines or blocks struck through, by strokes, loops or a wavy line; each is
% transcribed as \struck{} with one \note{} for the block. Red pencil frames
% formulas (5), (7), (17), (18), (19), (24), and marks a « ? » on page 54.
% Page 45 opens with the five lines that open page 44, on a narrower sheet;
% they are not repeated. The long sideways notes in the left margins of
% pages 43, 44, 54, 55 and 58 are recorded as \marginal{} where they can be
% read.
%
% Notation: he writes the Hilbert tensor product as ⊗ with a small « (2) »
% above the sign, and Λ with « n » and « (2) » above; both are set here as
% superscripts, \otimes^{(2)} and \Lambda^{n\,(2)}. The classes of
% Hilbert-Schmidt and Fredholm operators appear as L^{(2)}, L², L^{(1)}, L¹
% by turns; they are set uniformly as L^{(2)}(E,F) and L^{(1)}(E,F), with a
% note at the first occurrence. Ē is his conjugate space, E' the dual,
% (x,y) the scalar product, ρ_i the singular values of batch 2, S_n the
% symmetric group written as a gothic S.
%
% The dating is the inventory's own for the shelfmark. Its group,
% « MATHÉMATIQUES / RECHERCHE / Travaux », carries no date.
\documentclass[11pt,a4paper]{article}
% Le préambule commun à toutes les transcriptions du fonds Grothendieck.
%
% Il tient en une idée : **l'incertitude se compose, elle ne se résout pas.**
% Une transcription qui lisse un mot douteux a détruit la seule chose pour
% laquelle on la faisait — un lecteur doit pouvoir savoir, à chaque endroit,
% ce qui était sur le feuillet et ce qui a été ajouté. Les six macros qui
% suivent sont donc l'appareil critique, et non de la décoration.
%
% Les mêmes commandes sont reconnues par `scripts/render.mjs`, qui produit la
% vue de lecture du site. Une macro ajoutée ici doit l'être aussi là-bas,
% faute de quoi le rendu échouera bruyamment — ce qui est voulu : un rendu
% silencieusement faux est pire qu'un rendu absent.

\ProvidesPackage{grothendieck}[2026/08/08 Transcriptions du fonds Grothendieck]

\usepackage{amsmath,amssymb,amsthm}
\usepackage{mathrsfs}
\usepackage{stmaryrd}
% Les diagrammes commutatifs sont partout dans ce fonds : c'est la langue
% même de la géométrie algébrique de Grothendieck, et une transcription qui
% les rendrait en prose ne transcrirait rien.
\usepackage{tikz-cd}
% Français par défaut — c'est la langue des feuillets — mais l'anglais est
% chargé aussi : la traduction et le résumé sont des documents anglais, et la
% typographie française (espaces avant les deux-points) y serait une faute.
% Ces documents basculent par \selectlanguage{english} en tête de corps.
\usepackage[english,french]{babel}
\usepackage{fontspec}
\usepackage{geometry}
\geometry{a4paper,margin=25mm,marginparwidth=28mm}
\usepackage{marginnote}
\usepackage{xcolor}
% ulem, not soul. `soul`'s \st and \ul are built on a character-by-character
% scan that XeLaTeX's font machinery breaks: the document fails with an
% undefined control sequence rather than degrading. ulem does the same two jobs
% and is Unicode-engine safe. `normalem` stops it hijacking \emph, which the
% transcriptions use for Grothendieck's own emphasis.
\usepackage[normalem]{ulem}
\usepackage{hyperref}
\hypersetup{colorlinks=true,linkcolor=black,urlcolor=black,pdfborder={0 0 0}}

% --- Métadonnées du lot -----------------------------------------------------
% Reprises telles quelles par le script de rendu, qui les affiche en tête de
% la vue de lecture. Elles fixent ce que le fichier prétend couvrir : sans
% elles, une transcription flotte sans point d'attache dans un fonds de seize
% mille feuillets.

\newcommand{\folder}[1]{\gdef\grothFolder{#1}}
\newcommand{\batch}[1]{\gdef\grothBatch{#1}}
\newcommand{\leaves}[2]{\gdef\grothFirst{#1}\gdef\grothLast{#2}}
\newcommand{\foldertitle}[1]{\gdef\grothFoldertitle{#1}}
\gdef\grothFolder{?}\gdef\grothBatch{?}\gdef\grothFirst{?}\gdef\grothLast{?}\gdef\grothFoldertitle{}

% --- Le résumé --------------------------------------------------------------
%
% Il ouvre la lecture modernisée, sous le titre « Résumé », et il est composé
% en petit. Ce n'est pas un ornement : c'est le seul passage du document qui
% ne soit pas des mathématiques, et un lecteur doit pouvoir le reconnaître
% avant de l'avoir lu — puis le sauter s'il connaît déjà le sujet, ou s'y
% arrêter s'il ne le connaît pas. Le filet qui le suit marque la reprise du
% corps.

\newenvironment{resume}
  {\small\setlength{\parskip}{0.45em}}
  {\par\medskip\hrule\medskip}

% --- Mots-clés ---------------------------------------------------------------
%
% En anglais, en clôture du résumé de la lecture modernisée : le vocabulaire
% moderne sous lequel chercher ce que les feuillets construisent. Le manifeste
% du site (`npm run manifest`) les extrait pour étiqueter la cote entière —
% cette ligne est la source unique des tags, il n'y a pas de fichier à côté.
\newcommand{\keywords}[1]{\par\smallskip\noindent
  {\footnotesize\textsc{Keywords} --- \textit{#1}}\par}

% --- L'appareil critique ----------------------------------------------------

% Le numéro de feuillet, en marge. C'est l'ancre qui permet de régler un
% désaccord sur une formule en regardant *un* feuillet plutôt qu'une chemise
% de six cents. Le site s'en sert aussi pour tourner les pages du fac-similé
% quand on fait défiler la transcription.
\newcommand{\leaf}[1]{%
  \marginnote{\footnotesize\color{gray!70}\textsf{#1}}%
  \label{leaf:#1}\ignorespaces}

% Illisible : on ne devine pas. Un mot inventé qui se lit comme les autres est
% le pire résultat possible.
\newcommand{\ill}{\textcolor{red!60!black}{[\,\ldots\,]}}

% Lecture douteuse : le mot est donné, et signalé comme incertain.
\newcommand{\uncertain}[1]{\uline{#1}}

% Ajout de l'éditeur — une lettre manquante, un mot sous-entendu.
\newcommand{\add}[1]{\textcolor{blue!50!black}{[#1]}}

% Biffé par Grothendieck lui-même. Ce qu'il a rayé fait partie du document :
% une rature dit souvent où la pensée a changé de direction.
\newcommand{\struck}[1]{\sout{#1}}

% Note du transcripteur, sur l'état matériel du feuillet.
\newcommand{\note}[1]{\par\smallskip{\small\color{gray!60!black}\textit{[#1]}}\par\smallskip}

% Note marginale de Grothendieck, distincte du corps du texte.
\newcommand{\marginal}[1]{\par\smallskip\noindent\hspace{1em}%
  {\small\color{gray!60!black}#1}\par\smallskip}

% --- Titre ------------------------------------------------------------------

\AtBeginDocument{%
  \begin{center}
    {\large\bfseries Fonds Alexandre Grothendieck --- cote n\textsuperscript{o}~\grothFolder}\\[2pt]
    {\itshape\grothFoldertitle}\\[4pt]
    {\small feuillets \grothFirst--\grothLast\ (lot \grothBatch)}\\[2pt]
    {\footnotesize\color{gray!60!black}Transcription établie d'après le fac-similé numérisé
     par l'Université de Montpellier.\\
     Les lectures incertaines sont soulignées, les illisibles notées [\,\ldots\,],
     les ajouts entre crochets.}
  \end{center}
  \vspace{1em}}

\endinput


\folder{1}
\batch{3}
\pages{41}{60}
\dating{1953}
\watermark{Édition de démonstration}
\foldertitle{Analyse fonctionnelle : notes manuscrites (s.d.), lettre (1953)}

\begin{document}

\section{Caractérisation des opérateurs de Fredholm (fin)}

\note{Titre de l'éditeur. La page 41 achève le « II » des feuillets jaunis qui closent le lot 2 (pages 37 à 40) : après $\|u\|_1 = \sum \rho_i(u)$ et la définition de $L^{p}(E,F)$, le dual de $L_0(E,F)$.}

\page{41}
\emph{Dual de $L_0(E,F)$.} \struck{Il} Il \ill{} \uncertain{est} \ill{} \uncertain{identifié} \uncertain{isométriquement} avec $L^{1}(F,E)$, \struck{\ill{}} \uncertain{par}
\[
  (u, v) = \operatorname{Tr} vu
\]
\note{encadré} (\ill{} $E' \otimes F$ et $F' \otimes E$ \ill{} \ill{} \ill{}) et l'application \uncertain{naturelle} de $L^{1}(F,E)$ dans $(L_0(E,F))'$ \ill{} \uncertain{une} isométrie : si $v \in L^{1}(F,E)$ a \ill{} $\sum \rho_i\, \bar{f}_i \otimes e_i$ \note{le $\bar{f}_i$ est surchargé}, il suffit de prendre $u = \sum_{1}^{n} \rho_i\, \bar{e}_i \otimes f_i$ pour trouver $(u,v) = \sum_{1}^{n} \rho_i$. \note{ainsi sur la page, le facteur $\rho_i$ dans les deux sommes ; avec $v = \sum \rho_i \bar{f}_i \otimes e_i$ on attendrait $u = \sum \bar{e}_i \otimes f_i$} \uncertain{Montrons} \uncertain{qu'elle} \uncertain{est} \ill{}. En effet, si $\Phi$ est une forme linéaire continue sur $L_0(E,F)$, alors $\Phi(x' \otimes y)$ est une forme bilinéaire continue sur $E' \times F$, donc \uncertain{définit} une application linéaire $F \to E$, \uncertain{soit} $v$. \struck{On} $\langle u, \Phi\rangle = \operatorname{Tr} vu$, il \uncertain{suffit} de \uncertain{montrer} que $v$ est de Fredholm, \uncertain{et} \uncertain{alors} $|\operatorname{Tr} vu| \leq M\,\|u\|$ pour $u \in E' \otimes F$. \uncertain{Considérons} \struck{$p = vv^{*} \in L(E,E)$}, \struck{\uncertain{on} \uncertain{a} $\operatorname{Tr} Wv = \operatorname{Tr} vW$ ($u \in E' \otimes E$)}

\struck{$|\operatorname{Tr} uw| = |\operatorname{Tr} uvv^{*}| \leq |\operatorname{Tr} v^{*}uv|$ \ill{} \uncertain{donc} \ill{} $w$ \ill{} \uncertain{pourtant} \uncertain{que} $v$. \uncertain{On} \ill{}} \note{bloc de trois lignes biffé d'un trait ondulé}

$w = (vv^{*})^{1/2} \in L(E,E)$, \uncertain{donc} $v = Ww$, $w = Vv$, \uncertain{on} \uncertain{voit} aussitôt que $W$ \note{en interligne : « $= Vv$ », biffé} \uncertain{la} \uncertain{partie} \uncertain{unitaire} de $v$, \uncertain{montrons} que $w$ est \uncertain{opérateur} \uncertain{de} Fredholm, donc $v$ \note{en interligne : « $= Ww$ »} \ill{} (\uncertain{ce} qui \uncertain{ramène} \uncertain{à} : $F = E$). On \uncertain{montre} d'abord \uncertain{que} $w$ est \ill{}, \uncertain{par} \ill{} il \uncertain{contient} un \uncertain{sous-espace} $E_1 \subset E$ \uncertain{stable} \note{« stable » en interligne} \ill{} $w$, si $w$ \uncertain{n'était} \uncertain{pas} \uncertain{inversible}, ce qui \uncertain{amène} une \uncertain{contradiction}. On \uncertain{aura} \uncertain{donc} $w = \sum \rho_i\, \bar{f}_i \otimes e_i$, \uncertain{et} \uncertain{prenant} $u = \sum \bar{e}_i \otimes f_i$, on trouve $\sum \rho_i \leq M$. \uncertain{D'où} \uncertain{la} conclusion. --- \uncertain{Méthode} \note{« méthode » en interligne} \uncertain{plus} \uncertain{élémentaire} : \ill{} $\Phi$ une forme \uncertain{linéaire} \ill{} sur $L^{p}(E,F)$, \ill{} \uncertain{donnée} \uncertain{par} \uncertain{un} $v \in L^{q}(F,E)$, \ill{} \ill{} \ill{} \uncertain{opérateur} \ill{}. \ill{} \ill{} \ill{} \uncertain{tout} : \ill{}
\note{la page s'arrête sur cette ligne ; la page 42 est blanche}

\section{Produits tensoriels hilbertiens}

\note{Titre de sa main, en tête de la page 43, sous une ligne de plan « 1. Rappels et notations. Opérateurs de Fredholm » ; ce N° 2 court des pages 43 à 57. Il écrit le produit tensoriel hilbertien $\otimes$ avec un petit « (2) » au-dessus du signe ; ce « (2) » est rendu ici en exposant, $\otimes^{(2)}$. Les classes d'opérateurs de Hilbert-Schmidt et de Fredholm sont écrites tour à tour $L^{(2)}$, $L^{2}$, $L^{(1)}$, $L^{1}$ ; uniformisées en $L^{(2)}(E,F)$ et $L^{(1)}(E,F)$.}

\page{43}
\uncertain{1}. \emph{Rappels} \struck{\ill{}} \uncertain{et} \uncertain{notations}. \emph{Opérateurs de Fredholm} \note{les mots « Opérateurs de Fredholm » sont soulignés au crayon orange}

2. \emph{Produits tensoriels hilbertiens.} \struck{\uncertain{Des} Fredholm et \uncertain{des} Hilbert-Schmidt} \note{titre souligné deux fois ; en interligne au-dessous : « Tout le contenu de ce N° est \ill{} classique », avec « vraiment » ajouté au-dessus}

Soit $E$ un espace de Banach (sur le corps des \uncertain{nombres} complexes), désignons par $\bar{E}$ l'espace \struck{\ill{}} vectoriel obtenu en \uncertain{prenant} sur $E$ la structure vectorielle ayant la même structure additive que $E$, mais où la multiplication scalaire \struck{\ill{}} est donnée par $(\lambda, x) \to \bar{\lambda}.x$ \struck{$\bar{E}$ \ill{}} $\lambda \bar{x} = \overline{\bar{\lambda} x}$ \note{lecture incertaine des barres de conjugaison sur cette formule} On a bien là une structure d'espace vectoriel, \struck{et les \uncertain{fonctions} $\|x\|$ \ill{}} \struck{\ill{}}. L'application \uncertain{identique} de $E$ sur $\bar{E}$ est \uncertain{antilinéaire}, i.e. \uncertain{additive} et \uncertain{satisfait} à $\overline{\lambda x} = \bar{\lambda}\,\bar{x}$, \ill{} \ill{} \ill{} (\uncertain{en} \uncertain{particulier} \ill{} \ill{} \uncertain{de} Banach $\bar{E}$). \uncertain{Notons} \uncertain{aussi}
\[
  \text{(1)}\qquad \bar{\bar{E}} = E
\]
\note{au-dessous, une ligne biffée : « $\bar{E} = E$ \ill{} »}

\struck{\ill{} $F$ \ill{} \uncertain{application} \ill{} $E$ dans $F$ \ill{} $E$ dans $F$, $G$ \ill{} $E$, $F$, $G$ \uncertain{trois} \ill{} $(x,y) \to u(x,y)$ \ill{}} \note{bloc de quatre lignes biffé d'un trait ondulé et de traits obliques} bilinéaire de $E \times \bar{E}$ \struck{\ill{}} dans $G$ est dite \emph{sesquilinéaire} si \uncertain{elle} \uncertain{est} \uncertain{linéaire} \uncertain{en} $x$, antilinéaire en $y$ ; \uncertain{il} \uncertain{revient} \uncertain{au} \uncertain{même} \uncertain{de} \uncertain{dire} \uncertain{que} \uncertain{c'est} \uncertain{une} \uncertain{application} bilinéaire de $E \times \bar{F}$ dans $G$. Une application \uncertain{sesquilinéaire} \uncertain{correspond} \uncertain{donc} \uncertain{à} \uncertain{une} application linéaire de $E \otimes \bar{F}$ dans $G$. \note{en fin de ligne, à droite : « N.b. \uncertain{aussi} $L(E,\bar{F}) = \overline{L(E,F)}$ », lecture incertaine}

\struck{\ill{} $B$ \ill{} $\|x\|$ \ill{}} \note{deux lignes biffées de zigzags au bas de la page}

\marginal{en marge gauche, le long d'un trait oblique, sur plusieurs lignes : \ill{} $\bar{E}$ \ill{} $E'$ \ill{} $(x,y)$ \ill{}}

\page{44}
Notons que \uncertain{pour} \uncertain{deux} \uncertain{espaces} \uncertain{vectoriels} $E$, $F$, \uncertain{on} \uncertain{a} \uncertain{un} \uncertain{isomorphisme} \uncertain{canonique}
\[
  \text{(2)}\qquad \overline{E \otimes F} = \bar{E} \otimes \bar{F}
\]
En effet, l'application $(\bar{x}, \bar{y}) \to \overline{x \otimes y}$ de $\bar{E} \times \bar{F}$ \note{le signe entre $\bar{E}$ et $\bar{F}$ est surchargé} dans $\overline{E \otimes F}$ est \uncertain{manifestement} bilinéaire, donc se prolonge en une application linéaire $\bar{E} \otimes \bar{F} \to \overline{E \otimes F}$. \uncertain{On} \uncertain{vérifie} \uncertain{alors} \uncertain{aussitôt}, \uncertain{au} \uncertain{moyen} \uncertain{de} \uncertain{bases} \uncertain{dans} $E$ et $F$ par \uncertain{exemple}, \uncertain{que} \uncertain{c'est} \uncertain{un} \uncertain{isomorphisme} \uncertain{du} \uncertain{premier} \uncertain{espace} \uncertain{sur} le second.

\struck{Supposons} \uncertain{Rappelons}, \ill{} \ill{} \uncertain{préhilbertien} \ill{} \uncertain{un} \uncertain{espace} \uncertain{vectoriel} $E$, \uncertain{muni} \uncertain{d'une} \uncertain{forme} sesquilinéaire sur $E \times E$, \struck{\ill{} $x, y$} \uncertain{notée} $(x,y)$, \uncertain{satisfaisant} \uncertain{à} \uncertain{la} \uncertain{condition} $(x,x) \geq 0$ \uncertain{pour} \uncertain{tout} $x$. \uncertain{Alors} \struck{\ill{} $E \otimes$} $\|x\| = \sqrt{(x,x)}$ \uncertain{est} une \uncertain{semi-norme} sur $E$, et $E$ est \uncertain{dit} un espace de Hilbert si $E$ est \uncertain{séparé} et complet pour cette norme. \struck{Soit \ill{}} \struck{\ill{}} \uncertain{Soient} \ill{} \ill{} \uncertain{préhilbertiens} $E$ et $F$ \uncertain{deux} \uncertain{espaces} \struck{préhilbertiens} \ill{}. On \uncertain{peut} \uncertain{définir} \uncertain{sur} $E \otimes F$ une structure préhilbertienne \uncertain{d'une} \ill{} \ill{} : on \uncertain{doit} \uncertain{définir} une forme \uncertain{linéaire} \ill{} sur $(E \otimes F) \otimes \overline{E \otimes F}$ ; \uncertain{d'après} (2), \ill{} \ill{} \uncertain{isomorphe} \uncertain{canoniquement} \uncertain{à} $(E \otimes \bar{E}) \otimes (F \otimes \bar{F})$, \uncertain{il} \uncertain{faut} \uncertain{donc} \uncertain{définir} une forme linéaire \uncertain{sur} \ill{}, \uncertain{on} \uncertain{prend} \ill{}
\note{la phrase se poursuit page 47, après la reprise de la page 45}

\marginal{dans un cercle, en bas à gauche : « \uncertain{Introduire} $E \otimes F \subset L(\bar{E}, F)$ \uncertain{puis} $\operatorname{Tr}(v^{*}u)$ directement »}

\marginal{en marge gauche, sideways, le long de traits obliques, un brouillon de ce que la page 45 écrit au net : \ill{} $y \in E$ \ill{} $(x,y)$ \ill{} $E$ \ill{} $E'$ \ill{} ; (3) $E' = \bar{E}$ ($E$ \ill{}) ; $\|\bar{x}\| = $ \ill{}}

\page{45}
\note{les cinq premières lignes du feuillet, plus étroit, reprennent mot pour mot le début de la page 44, de « Notons que » à « $\bar{E} \otimes \bar{F} \to \overline{E \otimes F}$ » ; elles ne sont pas répétées}

Notons maintenant que si $E$ est un espace préhilbertien, alors la fonction $(\bar{x}, \bar{y}) \to \overline{(y,x)}$ sur $\bar{E} \times \bar{E}$ est une forme sesquilinéaire ; qui \uncertain{sera} \uncertain{notée} $(\bar{x}, \bar{y})$ ; \uncertain{qui} \struck{\ill{} $\bar{E}$ \ill{}} \uncertain{est} \struck{\ill{}} \uncertain{positive}, \uncertain{car} $(\bar{x}, \bar{x}) = (x, x)$. \uncertain{Elle} \uncertain{fait} \uncertain{de} \uncertain{l'ensemble} $\bar{E}$ \uncertain{un} \uncertain{espace} \uncertain{préhilbertien}, \uncertain{Comme} \ill{} \note{sous la ligne, en petit : « $(\bar{x},\bar{y}) = (y,x)$ », « $\|x\| = \|\bar{x}\|$ »} $\|x\| = \|\bar{x}\|$, \uncertain{la} \uncertain{norme} \uncertain{de} $\bar{E}$ \uncertain{coïncide} \uncertain{donc} \uncertain{avec} \uncertain{celle} \uncertain{de} $E$. \uncertain{Si} \uncertain{donc} $E$ est \uncertain{séparé}, $\bar{E}$ \uncertain{l'est}, et si $E$ est un espace de Hilbert, $\bar{E}$ est \uncertain{aussi} \uncertain{un} \uncertain{espace} \uncertain{de} \uncertain{Hilbert}. \uncertain{Dans} \uncertain{ce} \uncertain{dernier} \uncertain{cas}, \uncertain{on} \uncertain{sait} \uncertain{que} l'application \uncertain{qui} \uncertain{à} \struck{\ill{}} $y \in E$ \uncertain{fait} \uncertain{correspondre} la forme $x \to (x,y)$ sur $E$, est une \uncertain{application} \uncertain{antilinéaire} \note{« antilinéaire » en interligne} \uncertain{bijective} de $E$ \uncertain{sur} \uncertain{le} \uncertain{dual} $E'$, \uncertain{conservant} \uncertain{la} \uncertain{norme} ; \uncertain{par} \uncertain{suite} \ill{} \ill{} \uncertain{un} \uncertain{isomorphisme} \uncertain{canonique} \uncertain{d'espaces} \uncertain{normés} :
\[
  \text{(3)}\qquad E' = \bar{E} \qquad (E \text{ espace de Hilbert})
\]
\uncertain{Cet} \uncertain{isomorphisme} \uncertain{permet} \uncertain{de} \uncertain{considérer} $E'$ \uncertain{comme} \struck{\ill{}} \struck{\ill{}} \uncertain{un} \uncertain{espace} \uncertain{de} \uncertain{Hilbert}.
\note{le reste du feuillet est blanc ; la page 46 est blanche}

\page{47}
\uncertain{forme} \uncertain{quatre} \uncertain{fois} \uncertain{linéaire} sur $E \times \bar{E} \times F \times \bar{F}$. \struck{\uncertain{Pour} \ill{} \uncertain{ceci} \ill{}} \uncertain{Ceci} \uncertain{revient} \uncertain{par} \uncertain{suite} \uncertain{à} \uncertain{définir} \uncertain{dans} $E \otimes F$ \struck{une structure \uncertain{préhilbertienne}} \uncertain{une} \uncertain{forme} \uncertain{sesquilinéaire}, \struck{\ill{}} \uncertain{notée} $(u,v)$ \uncertain{donnée} \uncertain{par} :
\[
  \text{(5)}\qquad (x_1 \otimes y_1,\; x_2 \otimes y_2) = (x_1, x_2)\,(y_1, y_2)
\]
\note{encadré au crayon rouge ; aucun (4) n'apparaît ; dans le second membre, les virgules sont surchargées sur des signes $\otimes$} (\uncertain{car} \uncertain{le} \uncertain{deuxième} \uncertain{membre} \uncertain{de} (5) \uncertain{est} \uncertain{bien} \uncertain{linéaire} en $x_1 \in E$, $x_2 \in \bar{E}$, $y_1 \in F$, $y_2 \in \bar{F}$). \struck{\uncertain{Montrons} \uncertain{que} \ill{}} \uncertain{Montrons} \uncertain{que} \uncertain{cette} \uncertain{forme} \struck{définie positive} \uncertain{est} \uncertain{une} \uncertain{forme} \emph{définie positive},

\struck{i.e. que $(u,u) \geq 0$ \ill{} $(u,u) = 0 \Longleftrightarrow u = 0$. Soit $(e_i)_{i \in I}$ une base orthonormale dans $E$, $(f_j)_{j \in J}$ \ill{} \ill{} dans $F$, alors $(e_i \otimes f_j)_{(i,j) \in I \times J}$ \ill{}} \note{bloc de trois lignes biffé d'un trait ondulé}

En effet, si $u, v \in E \otimes F$, alors \struck{$u$ \ill{} \ill{}} \struck{$u, v \in E_0 \otimes F_0$, \ill{} \ill{}} \uncertain{où} $E_0$ et $F_0$ \struck{\ill{}} \uncertain{sont} \uncertain{deux} \uncertain{sous-espaces} \uncertain{de} \uncertain{dimension} \uncertain{finie} de $E$ \uncertain{et} \uncertain{de} $F$, \struck{Prenons alors une} \uncertain{soit} $(e_i)$ \note{en interligne} une base orthonormale \uncertain{dans} $E_0$, \uncertain{une} \uncertain{autre} $(f_j)$ \uncertain{dans} $F_0$, \struck{\ill{} $u = \sum$} \uncertain{alors} $(e_i \otimes f_j)$ \uncertain{est} \uncertain{une} \uncertain{base} \uncertain{de} \uncertain{l'espace} \uncertain{vectoriel} \ill{}, \uncertain{et} \uncertain{on} \uncertain{voit} \uncertain{sur} (5) \uncertain{que}
\[
  \text{(6)}\qquad (e_i \otimes f_j,\; e_{i'} \otimes f_{j'}) = \delta_{(i,j),(i',j')} \qquad (\text{indice de Kronecker}),
\]
\uncertain{ce} \uncertain{qui} \uncertain{prouve} \uncertain{que} $(u,v)$ induit \uncertain{sur} $E_0 \otimes F_0$ une forme \emph{définie positive}. \uncertain{Comme} $E \otimes F$ \uncertain{est} \uncertain{réunion} \uncertain{des} \uncertain{espaces} $E_0 \otimes F_0$, \uncertain{il} \uncertain{s'ensuit} \uncertain{que} $(u,v)$ \uncertain{est} \emph{définie} \uncertain{positive}.

\emph{Définition.} \uncertain{Soient} $E$, $F$ \uncertain{deux} \uncertain{espaces} de Hilbert. \uncertain{On} \uncertain{appelle} \uncertain{produit} \uncertain{tensoriel} \uncertain{hilbertien} de $E$ et $F$, \uncertain{et} \uncertain{on} \uncertain{note} $E \otimes^{(2)} F$, \uncertain{le} \uncertain{complété} \uncertain{de} \uncertain{l'espace} \struck{$E \otimes F$} \uncertain{préhilbertien}
\note{la phrase se poursuit page 48}

\page{48}
\uncertain{préhilbertien} $E \otimes F$, \struck{\ill{}} \uncertain{le} \uncertain{produit} \uncertain{scalaire} \ill{} \uncertain{défini} \uncertain{par} (5). \struck{\ill{} \uncertain{si} $u \in E \otimes^{(2)} F$ \ill{} \uncertain{on} \uncertain{note} $\|u\|_2$ \ill{}} \note{ligne serrée entre les lignes, biffée} $E \otimes^{(2)} F$ \uncertain{devient} \ill{} \uncertain{un} \uncertain{espace} \uncertain{de} Hilbert.

\emph{Prop 1.} \uncertain{Soient} $E$ et $F$ \uncertain{deux} \uncertain{espaces} de Hilbert. \struck{Si $x \in E$, $y \in F$, \ill{} $\|x \otimes y\|_2 = \|x\|\,\|y\|$. 2) Soit} \uncertain{Soit} $(e_i)_{i \in I}$ une base orthonormale \uncertain{dans} $E$, $(f_j)_{j \in J}$ une base orthonormale dans $F$, \uncertain{alors} $(e_i \otimes f_j)_{(i,j) \in I \times J}$ \uncertain{est} \uncertain{une} \uncertain{base} \uncertain{orthonormale} \uncertain{dans} $E \otimes^{(2)} F$. \note{l'énoncé est embrassé d'une accolade en marge gauche}

En effet, \uncertain{d'après} (6), \uncertain{le} \uncertain{système} $(e_i \otimes f_j)$ \uncertain{est} \uncertain{orthonormal}, \struck{\ill{} \ill{}} \struck{\ill{} \uncertain{dense} $E \otimes^{(2)} F$} \struck{\ill{}} \uncertain{et} \uncertain{il} \uncertain{est} \uncertain{total} : \uncertain{cela} \uncertain{résulte} \uncertain{de} \uncertain{ce} \uncertain{que} \uncertain{les} \uncertain{combinaisons} \uncertain{linéaires} \uncertain{des} $e_i \otimes f_j$ \uncertain{sont} \uncertain{denses} \uncertain{dans} $E \otimes^{(2)} F$. \uncertain{Pour} \uncertain{le} \uncertain{voir}, \uncertain{il} \uncertain{suffit} \uncertain{de} \uncertain{voir} \uncertain{que} \uncertain{tout} $u \in E \otimes F$ \uncertain{est} \uncertain{limite} \ill{} \ill{} \uncertain{si} $E_0$ \uncertain{est} \uncertain{engendré} \uncertain{par} \uncertain{des} \uncertain{combinaisons} \uncertain{finies} \uncertain{des} $e_i$, $F_0$ \ill{} \uncertain{des} $f_j$, \ill{} \uncertain{dans} $E$, $F_0$ \uncertain{dans} $F$, \uncertain{donc} \ill{} $E_0 \otimes F_0$ \uncertain{dans} \ill{} $E \otimes F$, \uncertain{donc} \ill{} \uncertain{dans} $E \otimes^{(2)} F$. \uncertain{Notons} \uncertain{aussi} \uncertain{deux} \uncertain{formules} :
\[
  \text{(7)}\qquad \|x \otimes y\|_2 = \|x\|\,\|y\|
\]
\note{encadré au crayon rouge} (\uncertain{qui} \uncertain{résulte} \uncertain{aussitôt} \uncertain{de} \struck{\ill{}} $(x \otimes y, x \otimes y) = (x,x)(y,y)$), \struck{\ill{}} \uncertain{et} \uncertain{aussi} \uncertain{que} l'isomorphisme \uncertain{canonique} $\overline{E \otimes F} = \bar{E} \otimes \bar{F}$ (\uncertain{cf.} (2)) \uncertain{conserve} \ill{} \uncertain{les} \uncertain{produits} \uncertain{scalaires}, \uncertain{donc} \uncertain{définit} \uncertain{un} \uncertain{isomorphisme} \uncertain{d'espaces} \uncertain{préhilbertiens}, \uncertain{d'où} \uncertain{un} \uncertain{isomorphisme} \uncertain{canonique}
\[
  \text{(8)}\qquad \overline{E \otimes^{(2)} F} = \bar{E} \otimes^{(2)} \bar{F}
\]

\page{49}
\note{les quatre premières lignes de la page sont enfermées dans une longue boucle partant du haut à gauche ; « Th 1 » et un premier « Exemple » y sont biffés}

\uncertain{L'exemple} \uncertain{le} \uncertain{plus} \uncertain{important} \uncertain{de} \uncertain{produit} \uncertain{tensoriel} \uncertain{hilbertien} \uncertain{est} \uncertain{donné} \uncertain{par} \uncertain{le} \struck{Th 1} \uncertain{Soit} $E$ \uncertain{un} \uncertain{espace} \uncertain{de} Hilbert, \struck{\ill{} $\mu$ \uncertain{une} \uncertain{mesure}} \uncertain{positive} \uncertain{sur} \uncertain{un} \uncertain{espace} \uncertain{localement} \uncertain{compact} $M$, \uncertain{considérons} \ill{} $L^{2}(\mu)$, \uncertain{et} \uncertain{l'espace} \uncertain{de} Hilbert

\struck{\emph{Exemple}} \struck{Soit} Soit $M$ \uncertain{un} \uncertain{espace} \uncertain{localement} \uncertain{compact}, $\mu$ \uncertain{une} \uncertain{mesure} \uncertain{positive} \uncertain{sur} $M$, \uncertain{considérons} \uncertain{l'espace} \uncertain{de} Hilbert $L^{2}(\mu)$ (\uncertain{où} $(\varphi, \psi) = \int \varphi\,\bar{\psi}\, d\mu$). \uncertain{Plus} \uncertain{généralement}, \uncertain{si} $E$ \uncertain{est} \uncertain{un} \uncertain{espace} \uncertain{de} Hilbert, $L^{2}_{E}(\mu)$ \uncertain{désigne} \uncertain{l'espace} \uncertain{des} \uncertain{fonctions} \ill{} \uncertain{sur} $M$ \uncertain{à} \uncertain{valeurs} \uncertain{dans} $E$, \uncertain{de} \uncertain{carré} \uncertain{intégrable}, \uncertain{muni} \uncertain{de} \uncertain{la} \uncertain{norme}
\[
  \text{(9)}\qquad \|f\|_2 = \Bigl(\int \struck{\ill{}}\, \|f(t)\|^{2}\, d\mu(t)\Bigr)^{1/2}
\]
\uncertain{qui} \uncertain{en} \uncertain{fait} \uncertain{un} \uncertain{espace} \uncertain{de} Banach (\ill{}) \uncertain{et} \uncertain{même} \uncertain{un} \uncertain{espace} \uncertain{de} Hilbert ; \uncertain{le} \uncertain{produit} \uncertain{scalaire} \uncertain{étant} \uncertain{défini} \uncertain{par} \uncertain{la} \uncertain{forme} sesquilinéaire
\[
  \text{(10)}\qquad (f, g) = \int \langle f(t), g(t)\rangle\, d\mu(t)
\]
\uncertain{dont} $L^{2}_{E}(\mu)$ \uncertain{est} \uncertain{un} \uncertain{espace} \uncertain{de} \uncertain{Hilbert}. \uncertain{On} \uncertain{a} \uncertain{une} \uncertain{application} \uncertain{bilinéaire} $(\varphi, a) \to \varphi.a$ \uncertain{de} $L^{2}(\mu) \times E$ \uncertain{dans} $L^{2}_{E}$, \uncertain{d'où} \uncertain{une} \uncertain{application} \uncertain{linéaire} $L^{2} \otimes E \to L^{2}_{E}$ \note{en interligne : « $u \to \tilde{u}$ »}, \uncertain{dont} \uncertain{l'image} \uncertain{est} \uncertain{dense} \uncertain{dans} $L^{2}_{E}$. \ill{} \uncertain{c'est} \uncertain{un} \uncertain{isomorphisme} \uncertain{préhilbertien} \ill{}, \uncertain{i.e.} $(\varphi.a, \psi.b) = (\varphi, \psi)(a, b)$, \uncertain{i.e.} $(\widetilde{\varphi \otimes a}, \widetilde{\psi \otimes b}) = (\varphi \otimes a, \psi \otimes b)$, \uncertain{d'où} $(\tilde{u}, \tilde{v}) = (u, v)$ \uncertain{pour} \uncertain{tout} $u, v \in L^{2} \otimes E$.

\page{50}
\uncertain{On} \uncertain{en} \uncertain{conclut} \uncertain{le}

\emph{Théorème 1.} \note{suivi d'un signe cerclé} \uncertain{Soit} $E$ \uncertain{un} \uncertain{espace} \uncertain{de} Hilbert, $\mu$ \uncertain{une} \uncertain{mesure} \uncertain{positive} \uncertain{sur} \uncertain{un} \uncertain{espace} \uncertain{localement} \uncertain{compact}, \uncertain{on} \uncertain{a} \uncertain{un} \uncertain{isomorphisme} \uncertain{canonique} \uncertain{d'espaces} \uncertain{hilbertiens}
\[
  \text{(11)}\qquad L^{2}(\mu) \otimes^{(2)} E = L^{2}_{E}(\mu)
\]
\note{l'énoncé est embrassé d'une accolade en marge gauche, comme les deux corollaires}

\emph{Corollaire 1.} $I$ \uncertain{étant} \uncertain{un} \uncertain{ensemble} \uncertain{quelconque} \uncertain{d'indices}, $E$ \uncertain{un} \uncertain{espace} \uncertain{de} Hilbert, \uncertain{on} \uncertain{a} \uncertain{un} \uncertain{isomorphisme} \uncertain{canonique}
\[
  \text{(12)}\qquad l^{2}(I) \otimes^{(2)} E = l^{2}_{E}(I)
\]
\uncertain{qui} \uncertain{à} \ill{} \uncertain{fait} \uncertain{correspondre} \struck{\ill{}} \uncertain{la} \uncertain{famille} $(a_i)_{i \in I}$ \uncertain{d'éléments} \uncertain{de} $E$, \struck{\ill{}} \ill{} \ill{} (\uncertain{avec} \struck{\ill{}} \uncertain{le} \uncertain{produit} \uncertain{scalaire} $((x_i), (y_i)) = \sum (x_i, y_i)$).

\emph{Corollaire 2.} \uncertain{Soient} $M$, $N$ \uncertain{deux} \uncertain{espaces} \uncertain{localement} \uncertain{compacts}, $\mu$, $\nu$ \uncertain{des} \uncertain{mesures} \uncertain{positives} \uncertain{dessus}, \uncertain{alors} \uncertain{on} \uncertain{a} \uncertain{un} \uncertain{isomorphisme} \uncertain{canonique} \uncertain{d'espaces} hilbertiens
\[
  \text{(13)}\qquad L^{2}(\mu) \otimes^{(2)} L^{2}(\nu) = L^{2}(\mu \otimes \nu)
\]
(\uncertain{où} $\mu \otimes \nu$ \uncertain{est} \uncertain{le} \uncertain{produit} \ill{} \uncertain{sur} $M \times N$). En effet, \uncertain{par} \uncertain{définition} $L^{2}(\mu \otimes \nu) = L^{2}_{L^{2}(\nu)}(\mu)$, \uncertain{et} \uncertain{on} \uncertain{applique} \uncertain{deux} \uncertain{fois} \uncertain{le} \uncertain{th.} 1. \uncertain{En} \uncertain{particulier}, $L^{2}(\mu)$ \uncertain{et} $L^{2}(\nu)$ \ill{} $l^{2}(I)$ \uncertain{et} $l^{2}(J)$, \ill{} \uncertain{retrouve} \uncertain{la} \uncertain{prop.} 1. \struck{\ill{}}

\uncertain{Soient} \uncertain{maintenant} $E$ et $F$ \uncertain{deux} \uncertain{espaces} de Hilbert, \ill{} \ill{} \uncertain{à} $E' = \bar{E}$, \struck{\ill{}} \uncertain{il} \uncertain{y} \uncertain{a} \uncertain{une} \uncertain{correspondance} \uncertain{entre} $\bar{E} \otimes F$ \uncertain{et} \uncertain{l'espace} \uncertain{des} \uncertain{applications} \uncertain{linéaires} \uncertain{de} \uncertain{rang} \uncertain{fini} \uncertain{de} $E$ \uncertain{dans} $F$, \uncertain{définie} \uncertain{par}
\[
  \text{(14)}\qquad (\bar{a} \otimes b).x = (x, a).b
\]
\note{suivi d'un point d'exclamation}

\page{51}
\struck{\uncertain{Montrons} \uncertain{que} \ill{}} \struck{\ill{}} \struck{(15) \ill{}} \struck{\uncertain{Comme} \ill{}} \note{quatre lignes biffées en tête de page} \ill{} (\uncertain{d'après} \uncertain{la} \uncertain{seconde} \struck{\ill{} (15) \uncertain{la} \uncertain{seconde}}) \ill{} \uncertain{l'espace} \uncertain{préhilbertien} \struck{\ill{}} $\bar{E} \otimes F$ \uncertain{s'identifie}, \uncertain{grâce} \uncertain{à} \uncertain{la} \uncertain{bijection} \uncertain{précédente}, \uncertain{et} $u = \bar{a} \otimes b$,
\[
  \text{(15)}\qquad (ux, y) = (u,\; \bar{x} \otimes y)
\]
\uncertain{d'où}, \uncertain{par} \uncertain{linéarité}, \uncertain{pour} \uncertain{tout} $u$ \ill{} \ill{} \uncertain{préhilbertien} $\bar{E} \otimes F$), \uncertain{formule} \uncertain{qui} \ill{} \uncertain{pour} \struck{\ill{}} $u \in \bar{E} \otimes F$. \uncertain{D'où} \uncertain{on} \uncertain{conclut} \uncertain{que}
\[
  |(ux, y)| \leq \|u\|_2\, \|\bar{x} \otimes y\|_2 = \|u\|_2\, \|\bar{x}\|\, \|y\| = \|u\|_2\, \|x\|\, \|y\|,
\]
\uncertain{d'où} \uncertain{résulte} \uncertain{que} \uncertain{l'opérateur} \ill{} \uncertain{norme} $\leq \|u\|_2$. \uncertain{Par} \uncertain{suite} $u \to \tilde{u}$ \uncertain{de} $\bar{E} \otimes F$ \uncertain{dans} $L(E,F)$ \uncertain{est} \uncertain{continue} \uncertain{et} \uncertain{se} \uncertain{prolonge} \uncertain{en} \uncertain{une} \uncertain{application} \uncertain{linéaire} \uncertain{continue} \uncertain{de} \struck{\ill{}} \uncertain{norme} $\leq 1$ \uncertain{de} $\bar{E} \otimes^{(2)} F$ \uncertain{dans} $L(E,F)$ \note{le signe après $\bar{E} \otimes^{(2)} F$, peut-être un $\subset$, n'est pas lisible}. \uncertain{Cette} \uncertain{application} \uncertain{est} \uncertain{d'ailleurs} \uncertain{biunivoque}, \uncertain{car} \uncertain{si} $(e_i)_{i \in I}$ \uncertain{est} \uncertain{une} \uncertain{base} \uncertain{orthonormale} \uncertain{dans} $E$ \note{en interligne : « $(\bar{e}_i)$ une base orthonormale dans $\bar{E}$ »}, $(f_j)_{j \in J}$ \uncertain{une} \uncertain{base} \uncertain{orthonormale} \uncertain{dans} $F$, \uncertain{alors} $\bar{E} \otimes^{(2)} F$ \uncertain{s'identifie} \uncertain{à} \uncertain{l'espace} \uncertain{des} \uncertain{matrices} \ill{} (\uncertain{prop.} 1) \ill{} \ill{} \uncertain{de} \uncertain{carré} \uncertain{sommable}, \uncertain{et} \uncertain{l'application} \ill{} \struck{\ill{}} \uncertain{à} $u = (u_{ij})_{(i,j) \in I \times J}$ \uncertain{fait} \uncertain{correspondre} \uncertain{l'opérateur} \struck{\ill{}} \uncertain{défini} \uncertain{par} \uncertain{la} \uncertain{matrice} $(u_{ij})$, \uncertain{donc} $u$ \uncertain{est} \uncertain{nul} \uncertain{si} \uncertain{cet} \uncertain{opérateur} \uncertain{est} \uncertain{nul}. \struck{Nous ne distinguerons pas par la notation \ill{} $\bar{E} \otimes^{(2)} F$ et les opérateurs qu'ils définissent}

\emph{Définition 2.} \uncertain{Soient} $E$ et $F$ \uncertain{deux} \uncertain{espaces} \uncertain{de} Hilbert. \uncertain{On} \uncertain{appelle} \uncertain{opérateur} \uncertain{de} Hilbert-Schmidt \uncertain{de} $E$ \uncertain{dans} $F$ \struck{\ill{}} \uncertain{tout} \uncertain{opérateur} \uncertain{linéaire} \uncertain{de} $E$ \uncertain{dans} $F$ \uncertain{défini} \uncertain{par} \uncertain{un} \uncertain{élément} \uncertain{de} $\bar{E} \otimes^{(2)} F$. \uncertain{L'ensemble} \uncertain{de} \uncertain{ces} \uncertain{opérateurs}, \uncertain{muni} \uncertain{de} \uncertain{la} \uncertain{structure}
\note{la phrase se poursuit page 52}

\page{52}
\uncertain{hilbertienne} \uncertain{définie} \uncertain{par} \uncertain{celle} \uncertain{de} $\bar{E} \otimes^{(2)} F$, \uncertain{sera} \uncertain{noté} $L^{(2)}(E,F)$ \note{en interligne : « $L^{(2)}(E)$ si $F = E$ »} ; \ill{} \ill{} \uncertain{de} Hilbert-Schmidt, \uncertain{on} \uncertain{écrira} \uncertain{si} $u \in L^{(2)}(E,F)$ \uncertain{sa} \uncertain{norme} $\|u\|_2$. \uncertain{On} \uncertain{a} \ill{} \uncertain{donc}
\[
  \text{(16)}\qquad \|u\|_{\infty} \leq \|u\|_2
\]
\uncertain{Les} \uncertain{opérateurs} \uncertain{de} Hilbert-Schmidt \uncertain{sont} \uncertain{des} \uncertain{applications} \uncertain{compactes} \struck{\ill{} \uncertain{dans} $L(E,F)$ \ill{} \uncertain{de} \ill{}}, \uncertain{c'est} \uncertain{même} \uncertain{un} \uncertain{opérateur} \ill{} \uncertain{compact} \uncertain{dans} \uncertain{les} \uncertain{opérateurs} $(u_{ij})$ \ill{}. \uncertain{Si} \ill{} \uncertain{bases} \uncertain{orthonormales} $(e_i)$, $(f_j)$, \uncertain{alors} \uncertain{la} \uncertain{matrice} \ill{} \uncertain{et} \uncertain{on} \uncertain{a} \uncertain{donc} $\|u\|_2 = \bigl(\sum |u_{ij}|^{2}\bigr)^{1/2}$ (\uncertain{matrices} \uncertain{de} Hilbert-Schmidt). \uncertain{Pour} \uncertain{les} \ill{} \uncertain{on} \uncertain{peut} \uncertain{identifier} $\bar{E} \otimes^{(2)} F$ \uncertain{et} $L^{(2)}(E,F)$.

\marginal{au crayon, en marge gauche, en oblique : « (\uncertain{mais} \uncertain{aussi} \ill{} $\|u\|$ \ill{} $\|\bar{u}\|$ \ill{}) »}

\emph{Théorème 2.} \uncertain{Soient} $E$ \uncertain{et} $F$ \uncertain{deux} \uncertain{espaces} \uncertain{de} Hilbert, \uncertain{soit} $u$ \uncertain{une} \struck{\ill{} \uncertain{H.-S.}} \uncertain{application} \uncertain{linéaire} \uncertain{continue} \uncertain{de} $E$ \uncertain{dans} $F$. \note{accolade en marge gauche}

1) \uncertain{Si} $u$ \uncertain{est} \uncertain{un} \uncertain{opérateur} \uncertain{de} H.-S., \struck{\ill{}} \uncertain{il} \uncertain{en} \uncertain{est} \uncertain{de} \uncertain{même} \uncertain{de} $Au$ \uncertain{lorsque} $A$ \note{le $A$ est surchargé} \uncertain{est} \uncertain{une} \uncertain{application} \uncertain{linéaire} \uncertain{continue} \uncertain{de} $F$ \uncertain{dans} \struck{\ill{}} \uncertain{un} \uncertain{espace} \uncertain{de} Hilbert $G$ (\uncertain{resp.} \uncertain{de} $uA$ \uncertain{lorsque} $A$ \ill{} \uncertain{dans} $E$) ; \uncertain{et} \uncertain{on} \uncertain{a}
\[
  \text{(17)}\qquad \|Au\|_2 \leq \|A\|\,\|u\|_2 \qquad (\text{resp. } \|uA\|_2 \leq \|A\|\,\|u\|_2)
\]
\note{encadré au crayon rouge ; un premier « (17) » est biffé au-dessous}

2) \uncertain{Pour} \uncertain{que} $u$ \uncertain{soit} \uncertain{un} \uncertain{opérateur} \uncertain{de} Hilbert-Schmidt, \uncertain{il} \uncertain{faut} \uncertain{et} \uncertain{il} \uncertain{suffit} \uncertain{que} $u^{*}$ \uncertain{soit} \uncertain{un} \uncertain{opérateur} \uncertain{de} H.-S. \uncertain{et} \uncertain{alors} $\|u\|_2 = \|u^{*}\|_2$ \note{encadré au crayon rouge}. \uncertain{Plus} \uncertain{précisément}, \uncertain{si} $F$ \uncertain{et} \ill{} $E$). (\uncertain{On} \uncertain{a}, \uncertain{pour} $u, v \in L^{(2)}(E,F)$
\[
  \text{(18)}\qquad \struck{(u, v) = (v^{*}, u^{*})}
\]
\note{ligne encadrée au crayon rouge puis biffée d'un trait ondulé ; suivie de « \uncertain{cor.} \ill{} \ill{} \uncertain{dans} $E$) »}

\emph{Th. 3.} \uncertain{Soit} $h = \sqrt{u^{*}u}$ \note{« $(u^{*}u)$ » en interligne} \ill{} \uncertain{Pour} \uncertain{que} $u$ \uncertain{soit} \uncertain{un} \uncertain{opérateur} \uncertain{de} H.-S., \uncertain{il} \uncertain{faut} \uncertain{et} \uncertain{il} \uncertain{suffit} \struck{\ill{} \uncertain{que} \ill{} \uncertain{soit} \ill{}} \note{deux lignes biffées, avec « (") » en interligne} \uncertain{que} \uncertain{les} \uncertain{valeurs} \uncertain{propres} $(\rho_i)$ \uncertain{de} $u$ \uncertain{soient} \uncertain{telles} \uncertain{que} $\sum \rho_i^{2} < +\infty$. \uncertain{Alors},
\[
  \text{(19)}\qquad \|u\|_2 = \|u^{*}\|_2 = \Bigl(\sum \rho_i^{2}\Bigr)^{1/2}
\]
\note{encadré au crayon rouge}

\struck{\ill{}} \emph{Corollaire 1.} $u$ \uncertain{H.-S.} $\Longleftrightarrow$ $u^{*}u$ \uncertain{Fredholm}, \uncertain{et} \uncertain{on} \uncertain{a} $\|u\|_2 = \sqrt{\operatorname{Tr} \struck{\ill{}}\, u^{*}u} = \sqrt{\|u^{*}u\|_1}$.

\marginal{en marge gauche, en oblique, l'amorce du même énoncé : « Corollaire », « $u$ H.-S. $\Longrightarrow$ », « $\|u\|_1$ » ; au crayon, au-dessous : « (\uncertain{Condition} \ill{}) »}

\page{53}
\emph{Démonstration.} \struck{Il suffit} 1) \uncertain{Prouvons} \uncertain{par} \uncertain{exemple} \uncertain{que} $Au$ \uncertain{est} \uncertain{un} \uncertain{opérateur} \uncertain{de} Hilbert-Schmidt \uncertain{si} $u$ \uncertain{l'est}, \uncertain{et} \uncertain{que} \struck{\ill{}} $\|Au\|_2 \leq \|A\|\,\|u\|_2$. \uncertain{En} \uncertain{effet}, \uncertain{soit} $(e_i)$ \uncertain{une} \uncertain{base} \uncertain{orthonormale} \uncertain{dans} $E$, \uncertain{on} \uncertain{a} \struck{$\sum \|Aue_i\|$} $\sum \|Aue_i\|^{2} \leq \struck{\ill{}} \sum \|A\|^{2} \struck{\ill{}} \|ue_i\|^{2} = \|A\|^{2} \sum \|ue_i\|^{2}$, \uncertain{d'où} \struck{\ill{}} \uncertain{résulte} \uncertain{aussitôt} \uncertain{l'assertion}. \ill{} \ill{} \ill{} \struck{$\sum \|ue_i\|^{2}$ \ill{}} \ill{} \ill{} \uncertain{(?)}. \struck{\ill{}} \note{plusieurs lignes biffées} 2) \struck{\ill{}} \ill{} \ill{} \ill{}. \uncertain{Considérons} \ill{} $\overline{\bar{E} \otimes^{(2)} F} = E \otimes^{(2)} \bar{F}$ \note{d'après (8)}, \ill{} \uncertain{isomorphisme} \ill{} \struck{\ill{} $= E \otimes^{(2)} \bar{F}$}, \uncertain{donc} \ill{} : $\bar{F} \otimes^{(2)} E$. \struck{\ill{}} \uncertain{un} \uncertain{isomorphisme} \ill{} \uncertain{de} \uncertain{Hilbert} \uncertain{de} \struck{$\bar{E} \otimes F$} $L^{(2)}(E,F)$ \uncertain{sur} $L^{(2)}(F,E)$, \uncertain{et} \uncertain{on} \uncertain{vérifie} \uncertain{aussitôt} \uncertain{que} \uncertain{cette} \uncertain{application} \uncertain{est} \uncertain{donnée} \uncertain{par} $u \to u^{*}$ ($u \in L^{(2)}(E,F)$), \uncertain{d'où} 2).

3) \uncertain{On} \uncertain{sait} \uncertain{qu'il} \uncertain{existe} \ill{} \ill{} $u = Uh$, \uncertain{où} $U$ \uncertain{est} \uncertain{une} \uncertain{isométrie} \uncertain{partielle} \uncertain{de} $E$ \uncertain{dans} $F$, \uncertain{et} $h = Vu$, \uncertain{où} $V$ \uncertain{est} \uncertain{une} \uncertain{application} \uncertain{partiellement} \uncertain{isométrique} \uncertain{de} $F$ \uncertain{dans} $E$. \uncertain{On} \uncertain{en} \uncertain{conclut}, \uncertain{par} 1), \uncertain{que} \struck{\ill{}} \uncertain{si} $u$ \uncertain{est} \uncertain{un} \uncertain{opérateur} \uncertain{de} H.-S., \uncertain{il} \uncertain{en} \uncertain{est} \uncertain{de} \uncertain{même} \uncertain{de} $h$, \uncertain{et} \uncertain{inversement} \struck{\ill{}} \uncertain{si} $h$ \uncertain{est} \uncertain{de} H.-S., \uncertain{et} \uncertain{que} \struck{\ill{}} \uncertain{alors} $\|u\|_2 = \|h\|_2$. \uncertain{D'autre} \uncertain{part}, \uncertain{pour} \uncertain{que} $h$ \uncertain{soit} \uncertain{un} \uncertain{opérateur} \uncertain{de} Hilbert-Schmidt, \uncertain{il} \uncertain{faut} \uncertain{et} \uncertain{il} \uncertain{suffit} \uncertain{que} \ill{}. \uncertain{Alors} \uncertain{on} \uncertain{sait}
\note{la phrase se poursuit page 54}

\page{54}
\uncertain{que} \uncertain{la} \ill{} \uncertain{une} \uncertain{base} \uncertain{orthonormale} \ill{} \uncertain{de} $E$, \uncertain{la} \ill{} \uncertain{dans} \ill{} \uncertain{avec} \uncertain{la} \uncertain{matrice} \uncertain{diagonale} \ill{}, \struck{\ill{}} \uncertain{alors} \uncertain{la} \uncertain{matrice} \uncertain{diagonale} \ill{} \uncertain{les} \uncertain{valeurs} \uncertain{propres} $(\rho_i)$ \uncertain{de} $h$, \uncertain{de} \ill{}, \uncertain{le} \ill{} \uncertain{de} H.-S. \uncertain{si} \uncertain{et} \uncertain{seulement} \uncertain{si} \ill{} \ill{} \uncertain{est} \uncertain{de} \uncertain{carré} \uncertain{sommable}, \uncertain{i.e.} \uncertain{si} $\bigl(\sum |\rho_i|^{2}\bigr)^{1/2} < +\infty$, \uncertain{et} \uncertain{alors} $\|h\|_2 = \bigl(\sum \rho_i^{2}\bigr)^{1/2}$. \uncertain{Le} \uncertain{th.} 3 \uncertain{est} \uncertain{donc} \uncertain{démontré}.

\uncertain{Les} \uncertain{relations} \uncertain{entre} \uncertain{opérateurs} \uncertain{de} H.-S. \uncertain{et} \uncertain{opérateurs} \uncertain{de} Fredholm \uncertain{sont} \uncertain{exprimées} \uncertain{dans} \uncertain{le}

\emph{Th 3.} \note{ce théorème porte le même numéro que le Th. 3 de la page 52 ; le « 3 » est surchargé} \uncertain{Soient} $E$ et $F$ \uncertain{deux} \uncertain{espaces} \uncertain{de} Hilbert. \uncertain{L'application} \uncertain{linéaire} \uncertain{canonique} \uncertain{de} $\bar{E} \otimes^{(2)} F$ \uncertain{dans} $L(E,F)$ \uncertain{est} \uncertain{biunivoque}. \uncertain{Pour} \uncertain{qu'une} \uncertain{application} \uncertain{linéaire} $u$ \uncertain{de} $E$ \uncertain{dans} $F$ \uncertain{soit} \uncertain{application} \uncertain{de} Fredholm, \uncertain{il} \uncertain{faut} \uncertain{et} \uncertain{il} \uncertain{suffit} \uncertain{que} $h = \sqrt{u^{*}u}$ \uncertain{le} \uncertain{soit} \struck{\uncertain{Pour}} \uncertain{Pour} \uncertain{ceci}, \uncertain{il} \uncertain{faut} \uncertain{et} \uncertain{il} \uncertain{suffit} \uncertain{que} \uncertain{les} \uncertain{valeurs} \uncertain{propres} $(\rho_i)$ \ill{} \struck{\ill{} \uncertain{soit} \uncertain{compact} \uncertain{et} \uncertain{que} \ill{} \uncertain{les} \uncertain{valeurs} \uncertain{propres} \ill{} \uncertain{de} $h$ \uncertain{soient} \uncertain{sommables}, \uncertain{et} \uncertain{alors} \ill{} : $\sum \rho_i < +\infty$, \uncertain{et} \uncertain{alors}} \note{deux lignes biffées, la première tachée de rouge ; accolade en marge gauche sur l'énoncé} \struck{\ill{}}
\[
  \text{(20)}\qquad \|u\|_1 = \|h\|_1 = \sum \rho_i
\]
\uncertain{On} \uncertain{peut} \uncertain{aussi} \uncertain{y} \uncertain{remplacer} $h$ \uncertain{par} $\sqrt{uu^{*}}$. \struck{\ill{}}

\uncertain{Que} $u$ \uncertain{et} $h$ \uncertain{soient} \uncertain{simultanément} \uncertain{de} Fredholm, \uncertain{et} \uncertain{que} $\|u\|_1 = \|h\|_1$, \uncertain{résulte} \uncertain{du} \uncertain{fait} \uncertain{qu'on} \uncertain{a} $u = Uh$, $h = Vu$, \uncertain{avec} $\|U\| \leq 1$, $\|V\| \leq 1$. \struck{\ill{}} \uncertain{Alors}, $h$ \uncertain{sera} \uncertain{un} \uncertain{opérateur} \uncertain{compact}, \struck{\uncertain{s'il} \uncertain{en} \uncertain{est} \ill{}} \uncertain{donc} \struck{$h = \sum \lambda_i\, \bar{a}_i \otimes a_i$} \uncertain{de} \uncertain{la} \uncertain{forme} $\sum \rho_i\, \bar{a}_i \otimes a_i$ \note{le $\rho$ corrigé sur un $\lambda$}, \uncertain{où} \uncertain{les} $(a_i)$ \uncertain{forment} \uncertain{un} \uncertain{système} \uncertain{orthonormal}, \uncertain{les} $\rho_i$ \uncertain{sont} \uncertain{les} \uncertain{valeurs} \uncertain{propres} \uncertain{de} $h$. \uncertain{Soit} $p_n$ \uncertain{le} \uncertain{projecteur} \uncertain{de} $E$ \uncertain{sur} \uncertain{l'espace} \uncertain{engendré} \uncertain{par} \uncertain{les} $a_1, \ldots, a_n$, \uncertain{on} \uncertain{aura} \uncertain{donc}
\[
  \sum_{1}^{n} \rho_i = \operatorname{Tr} p_n h \leq \|p_n h\|_1 \leq \|h\|_1,
\]
\uncertain{d'où} $\sum \rho_i \leq \|h\|_1 \leq \|u\|_1$. \uncertain{Réciproquement}, \uncertain{supposons}
\note{la phrase se poursuit page 55}

\marginal{en marge gauche, sideways, reliée par une boucle aux deux lignes biffées de l'énoncé : « \ill{} \uncertain{dans} \uncertain{2} \uncertain{premières} \ill{} ; \ill{} \uncertain{des} \uncertain{théorèmes} \ill{} ; \uncertain{d'où} \ill{} \uncertain{signalé} \uncertain{de} \uncertain{la} \uncertain{fin} \uncertain{des} ; § 1, N° 1, (\uncertain{petites} \uncertain{corrections}) » ; au crayon rouge, un « ? »}

\page{55}
\ill{} \uncertain{le} \ill{} \uncertain{et} \uncertain{la} \uncertain{série} \uncertain{de} \uncertain{ses} \uncertain{valeurs} \uncertain{propres} \note{« $(\rho_i)$ » en interligne} \ill{} $h = \sum \rho_i\, \bar{a}_i \otimes a_i$, \uncertain{où} $(a_i)$ \ill{} \uncertain{orthonormal}, \uncertain{d'où} $\|h\|_1 = \sum \rho_i\, \|a_i\|^{2} = \sum \rho_i$, \uncertain{ce} \uncertain{qui} \uncertain{prouve} \uncertain{que} $h$ \uncertain{est} \uncertain{de} \uncertain{Fredholm}, \uncertain{et} \uncertain{on} \uncertain{a} (20). \uncertain{Le} \uncertain{cas} \uncertain{où} \uncertain{l'on} \uncertain{remplace} $h$ \uncertain{par} $h' = \sqrt{uu^{*}}$, \uncertain{on} \uncertain{peut} \ill{} \uncertain{soit} \uncertain{directement}, \uncertain{soit} \uncertain{par} \uncertain{application} \uncertain{du} \uncertain{fait} \uncertain{que} \struck{\ill{}} \struck{\ill{}} \note{en interligne : « $L(E,F)$ »} \ill{} $\sqrt{u^{*}u}$ \uncertain{et} $\sqrt{uu^{*}}$ \uncertain{ont} \uncertain{les} \uncertain{mêmes} \uncertain{valeurs} \uncertain{propres} \note{« les mêmes valeurs propres » est biffé puis récrit}. \uncertain{En} \uncertain{effet} \uncertain{si} $(\lambda_i)$ \ill{} $h = \sqrt{u^{*}u}$, \uncertain{on} \uncertain{aura} $h = \sum \rho_i\, \bar{a}_i \otimes a_i$ \ill{} \uncertain{isométrique} \uncertain{sur} \uncertain{l'espace} \uncertain{engendré} \ill{}
\[
  u = \sum \rho_i\, \bar{a}_i \otimes b_i,
\]
\uncertain{où} $(b_i) = (Ua_i)$ \uncertain{orthonormal} \note{« dans $F$ » en interligne} \ill{} \uncertain{système} \uncertain{orthonormal}, \uncertain{d'où}
\[
  u^{*} = \sum \rho_i\, \bar{b}_i \otimes a_i, \quad \text{et}\quad uu^{*} = \sum_{i,j} \rho_i \rho_j\, (a_i, a_j)\, \bar{b}_j \otimes b_i = \sum \rho_i^{2}\, \bar{b}_i \otimes b_i,
\]
\[
  \text{d'où}\quad \sqrt{uu^{*}} = \sum \rho_i\, \bar{b}_i \otimes b_i,
\]

\emph{Corollaire 1.} \uncertain{Les} \uncertain{opérateurs} \uncertain{de} Fredholm \uncertain{de} $E$ \uncertain{dans} $F$ \uncertain{sont} \uncertain{exactement} \uncertain{les} \struck{\ill{}} \uncertain{opérateurs} \uncertain{qui} \uncertain{peuvent} \uncertain{se} \uncertain{mettre} \uncertain{sous} \uncertain{la} \uncertain{forme}
\[
  \text{(21)}\qquad u = \sum \rho_i\, \bar{a}_i \otimes b_i
\]
\note{le numéro est cerclé, la formule encadrée} \uncertain{où} $(a_i)$ \uncertain{et} $(b_i)$ \uncertain{sont} \uncertain{des} \uncertain{systèmes} \uncertain{orthonormaux} \uncertain{dans} $E$ (\uncertain{resp.} \uncertain{dans} $F$), \uncertain{et} $(\rho_i)$ \uncertain{une} \uncertain{suite} \uncertain{de} \uncertain{nombres} $\geq 0$ \uncertain{sommable} ; \uncertain{on} \uncertain{a} \uncertain{alors} $\|u\|_1 = \sum \rho_i$. \uncertain{Comme} $\bar{E} \otimes^{(2)} E = L(E)$ \ill{} \uncertain{biunivoque}, \ill{} \uncertain{par} \uncertain{le} \uncertain{produit} \ill{} \uncertain{opérateur} \uncertain{de} Fredholm \uncertain{dans} \ill{} \uncertain{l'espace} \uncertain{de} Hilbert \ill{} ; \struck{\ill{} 2. \uncertain{Tr} \uncertain{de} \ill{}} \struck{\ill{}} \uncertain{la} \uncertain{formule} (20) \uncertain{peut} \uncertain{s'écrire} \uncertain{aussi} \note{« s'écrire aussi » en interligne}
\[
  \text{(22)}\qquad \|u\|_1 = \operatorname{Tr} \sqrt{u^{*}u} = \operatorname{Tr} \sqrt{uu^{*}}
\]
\note{le numéro est cerclé, la formule encadrée}

\emph{Corollaire 2.} \uncertain{Si} $u$ \uncertain{est} \uncertain{un} \uncertain{opérateur} \uncertain{de} Fredholm \uncertain{de} $E$ \uncertain{dans} $F$,
\[
  \text{(23)}\qquad \|u\|_1 = \sup_{\|A\| \leq 1} |\operatorname{Tr} Au| \qquad (A \in L(F,E))
\]
\note{le numéro est cerclé ; « $\|A\| \leq 1$ » est ajouté en interligne} (\uncertain{car} \ill{}, \uncertain{le} \uncertain{dual} \uncertain{de} $E' \otimes^{(2)} F$ \ill{} $= B(E', F) \approx L(F,E)$) \note{coche au crayon rouge} \struck{\ill{}} \uncertain{les} \uncertain{opérateurs} \uncertain{de} Fredholm \uncertain{et} \uncertain{opérateurs} \uncertain{de} H.-S. \ill{} \ill{} \uncertain{dans} \uncertain{l'énoncé} \ill{}

\emph{Th. 4.} \uncertain{Soient} $E$, $F$, $G$ \uncertain{trois} \uncertain{espaces} \uncertain{de} Hilbert, $u \in L^{(2)}(E,F)$, $v \in L^{(2)}(F,G)$. \uncertain{Alors} $vu \in L^{(1)}(E,G)$, \uncertain{et} \struck{\ill{}}
\[
  \text{(24)}\qquad \|vu\|_1 \leq \|v\|_2\, \|u\|_2
\]
\note{le numéro est cerclé, la formule encadrée au crayon rouge}

\marginal{en marge gauche, sideways, sur une dizaine de lignes, des corollaires en brouillon : « \uncertain{Utiliser} \uncertain{aussi} \ill{} ; \uncertain{Corollaire} 3 \ill{} \uncertain{par} \uncertain{le} \uncertain{th.} 2, \uncertain{puis} (19), \ill{} ; \uncertain{pour} \uncertain{que} $u$ \uncertain{soit} \ill{} \uncertain{que} $u^{*}u$ \uncertain{soit} \uncertain{un} \uncertain{opérateur} \uncertain{de} Fredholm, \uncertain{et} \uncertain{alors} (23) $\|u\|_2 = \sqrt{\operatorname{Tr}(u^{*}u)}$ ; \ill{} $= \sqrt{\|u^{*}u\|_1}$ ; \uncertain{Corollaire} \ill{} $L^{(1)}(E,F) = \bar{E} \otimes F \subset E' \otimes^{(2)} F$ \ill{} $L^{(2)}(E,F)$ \ill{} $L^{(1)}(E)$ \ill{} » ; le (23) de la marge n'est pas celui du corps de la page}

\page{56}
\uncertain{De} \uncertain{plus}, \uncertain{si} $u, v \in L^{(2)}(E,F)$, \uncertain{on} \uncertain{a}
\[
  \text{(25)}\qquad (u, v) = \operatorname{Tr} uv^{*} = \operatorname{Tr} v^{*}u
\]
\note{le numéro est cerclé, la formule encadrée}

\emph{Démonstration.} \uncertain{Pour} \uncertain{prouver} (24), \uncertain{il} \uncertain{suffit} \uncertain{par} \uncertain{raison} \uncertain{de} \uncertain{continuité} \struck{\ill{}} \uncertain{de} \uncertain{le} \uncertain{prouver} \uncertain{au} \uncertain{cas} \uncertain{où} $u, v$ \uncertain{sont} \uncertain{de} \uncertain{rang} \uncertain{fini} (\uncertain{alors} \uncertain{l'application} \uncertain{bilinéaire} $(u,v) \to vu$ \ill{} \uncertain{par} \uncertain{continuité} \uncertain{en} \uncertain{une} \uncertain{application} \uncertain{bilinéaire} \uncertain{de} \uncertain{norme} $\leq 1$ \uncertain{de} $L^{(2)}(E,F) \times L^{(2)}(F,G)$ \uncertain{dans} $L^{(1)}(E,G)$, \uncertain{qui} \uncertain{doit} \uncertain{bien} \struck{\ill{}} \uncertain{ne} \uncertain{pouvoir} \uncertain{être} \uncertain{que} $(u,v) \to vu$). \uncertain{Or}, \struck{on} \uncertain{d'après} (22), \uncertain{on} \uncertain{a} \ill{} $|\operatorname{Tr} Avu| \leq \|Av\|_2\, \|u\|_2$ \uncertain{pour} $A \in L(G,E)$, $\|A\| \leq 1$. \uncertain{Comme} $\|Av\|_2 \leq \|A\|\,\|v\|_2$ (\uncertain{formule} (17)), \uncertain{on} \uncertain{en} \uncertain{conclut}, \uncertain{prenant} \ill{} \uncertain{le} \uncertain{sup} \uncertain{en} $A$ \note{en interligne : « posons $w = v^{*}$ »}, \ill{} ($v \in \struck{\bar{E} \otimes F}$) \note{en interligne, cerné : « ($v \in E' \otimes F$) »}, $\|v\| = \|v^{*}\|_2$, \uncertain{d'où} \struck{\ill{} $v^{*}$ \ill{}} $|\operatorname{Tr} v^{*}u| \leq \|v^{*}\|_2\, \|u\|_2$. \uncertain{Mais} \ill{} \uncertain{résulte} \uncertain{de} (25), \uncertain{puisque} $|(u,v)| \leq \|u\|_2\, \|v\|_2$. \uncertain{Reste} \uncertain{à} \uncertain{prouver} (25) ; \uncertain{par} \uncertain{raison} \uncertain{de} \uncertain{continuité}, \uncertain{il} \uncertain{suffit} \uncertain{de} \uncertain{le} \uncertain{faire} \uncertain{lorsque} $u, v$ \uncertain{sont} \uncertain{de} \uncertain{rang} \uncertain{fini}, \uncertain{et} \uncertain{même} \uncertain{lorsque} \struck{\ill{}} \struck{Ou} \uncertain{Or} \uncertain{le} \uncertain{prouver} \uncertain{lorsque} \ill{} \uncertain{fini} \uncertain{est} \uncertain{trivial} : \ill{} \uncertain{lorsque} \ill{} \uncertain{est}
\[
  u = \bar{a} \otimes b, \quad v = \bar{a}' \otimes b', \quad \text{d'où}\quad v^{*} = \bar{b}' \otimes a',
\]
\[
  v^{*}u = (b, b')\, \bar{a} \otimes a', \qquad \struck{uv^{*} = (a', a)\, \bar{b} \otimes \ill{}} \qquad \text{d'où}
\]
\[
  (u, v) = (\bar{a}, \bar{a}')\,(b, b') = \struck{(a^{*}, a')}\,(b, b'), \qquad \operatorname{Tr} v^{*}u = (a', a)\,(b, b').
\]
\note{au-dessus du facteur biffé, en interligne : « $(a', a)$ »}
\uncertain{Comme} $\operatorname{Tr} v^{*}u = \operatorname{Tr} uv^{*}$ (\struck{\ill{}} \note{en interligne : « \uncertain{formule} \uncertain{de} »} \ill{} \uncertain{pour} \uncertain{2} \uncertain{traces}) \uncertain{lorsque} $u, v$ \uncertain{de} \uncertain{rang} \uncertain{fini}. \uncertain{Remarquons} \uncertain{d'ailleurs} \uncertain{que} (25) \uncertain{est} \ill{} \uncertain{par} \ill{} \uncertain{dans} \uncertain{le} \uncertain{théorème}, \uncertain{donc} \uncertain{aussi} : \uncertain{la} \uncertain{première} \uncertain{partie} \uncertain{du} \uncertain{th.} \uncertain{est} \uncertain{démontrée}, \uncertain{et} \uncertain{la} \uncertain{formule} (25), \uncertain{valable} \uncertain{pour} $u \in L^{(2)}(E,F)$, $v \in L^{(2)}(E,F)$, \uncertain{peut} \uncertain{s'écrire} \uncertain{par} \uncertain{continuité} \ill{}

\page{57}
\emph{Corollaire 1.} \uncertain{Pour} \uncertain{que} $u \in L(E,F)$ \uncertain{soit} \struck{\ill{}} \uncertain{de} H.-S., \uncertain{il} \uncertain{faut} \uncertain{et} \uncertain{il} \uncertain{suffit} \uncertain{que} \struck{\ill{}} $vu$ \uncertain{soit} \uncertain{opérateur} \uncertain{de} Fredholm \uncertain{pour} \struck{\ill{}} \uncertain{tout} $v \in L^{(2)}(F,E)$. \uncertain{Alors} \struck{\uncertain{Alors}}
\[
  \text{(26)}\qquad \|u\|_2 = \sup_{\|v\|_2 \leq 1} |\operatorname{Tr} vu| \qquad (v \in L^{(2)}(F,E))
\]
\uncertain{En} \uncertain{effet}, \ill{} \uncertain{de} \uncertain{la} \uncertain{prop.} \ill{} \uncertain{la} \uncertain{condition} \uncertain{est} \uncertain{nécessaire}, \uncertain{et} \uncertain{d'autre} \uncertain{part} \uncertain{la} \uncertain{formule} (25) \uncertain{s'écrit}
\[
  \|u\|_2 = \sup_{\|w\|_2 \leq 1} |\operatorname{Tr} w^{*}u| \qquad (w \in L^{(2)}(E,F))
\]
\uncertain{ce} \uncertain{qui} \uncertain{résulte} \uncertain{immédiatement} \uncertain{de} $\operatorname{Tr} w^{*}u = (u, w)$. \uncertain{Réciproquement}, \uncertain{supposons} \uncertain{que} $vu \in L^{(1)}(E,E)$ \uncertain{pour} \struck{\ill{}} \uncertain{tout} $v \in L^{(2)}(F,E)$, \uncertain{l'application} $v \to vu$ \uncertain{de} $L^{(2)}(F,E)$ \uncertain{dans} $L^{(1)}(E)$ \note{un signe biffé après $L^{(1)}(E)$, ici et deux lignes plus bas} \uncertain{est} \uncertain{continue} ; \uncertain{car} \uncertain{elle} \uncertain{est} \ill{} \uncertain{fermée} (\uncertain{comme} \uncertain{elle} \uncertain{est} \uncertain{continue} \uncertain{pour} \uncertain{la} \uncertain{topologie} \uncertain{induite} \uncertain{par} $L^{(1)}(E)$ \ill{} \uncertain{la} \uncertain{topologie} \uncertain{de} \uncertain{la} \uncertain{norme} \uncertain{plus} \uncertain{faible} \uncertain{dans} $L(E)$, \ill{} \uncertain{continue}, \uncertain{donc} \uncertain{par} \uncertain{suite} \ill{} \ill{} \ill{}) \ill{} \uncertain{par} \uncertain{suite} \uncertain{la} \uncertain{forme} \struck{\ill{}} $w \to \operatorname{Tr} w^{*}u$ \note{le $w^{*}$ est corrigé} \uncertain{sur} $\bar{E} \otimes^{(2)} F$ \uncertain{est} \uncertain{continue}, \uncertain{donc} \uncertain{définie} \uncertain{par} \uncertain{un} \uncertain{élément} $u_0$ \uncertain{de} $\bar{E} \otimes^{(2)} F$ \struck{\ill{}} \ill{} : $\operatorname{Tr}(w^{*}u) = (u_0, w)$ ; \uncertain{prenant} $w = \bar{x} \otimes y$, \uncertain{on} \uncertain{trouve} \ill{} $(ux, y) = (u_0 x, y)$ \uncertain{pour} \uncertain{tout} $x \in E$, $y \in F$, \uncertain{d'où} $u = u_0$ \ill{}, \uncertain{d'où} \uncertain{le} \uncertain{corollaire} \ill{}

\emph{Corollaire 2.} \uncertain{Pour} \uncertain{que} $u \in L(E)$ \uncertain{soit} \uncertain{opérateur} \uncertain{de} Fredholm, \uncertain{il} \uncertain{faut} \uncertain{et} \uncertain{il} \uncertain{suffit} \uncertain{qu'il} \uncertain{soit} \uncertain{de} \uncertain{la} \uncertain{forme} $u = vw$, \uncertain{où} $u, v \in L^{(2)}(E)$ \note{sic : $u, v$ sur la page, là où l'on attend $v, w$}, \uncertain{et} \uncertain{on} \uncertain{peut} \uncertain{choisir} $\|v\|_2 = \|w\|_2 = \sqrt{\|u\|_1}$ \note{les lettres $u$, $v$, $w$ se confondent dans cette ligne ; lecture d'après le sens ; accolade en marge gauche sur l'énoncé}. \uncertain{Il} \uncertain{suffit} \uncertain{d'écrire}, \uncertain{si} $u \in L(E)$ \ill{}, \uncertain{résulte} \uncertain{du} \uncertain{th.} 4 \uncertain{la} \struck{\ill{}} \uncertain{suffisance}, \uncertain{et} \uncertain{réciproquement}, \uncertain{écrivons} $u = Uh = UW^{2}$, \uncertain{où} $W = \sqrt{h}$, \uncertain{on} \uncertain{a} (\uncertain{nécessité}) $\|W\|_2^{2} = \|h\|_1 = \|u\|_1$, \uncertain{et} \uncertain{prenons} $v = UW$, \uncertain{on} \uncertain{a} $\|v\|_2 \leq \|W\|_2$, \uncertain{et} \uncertain{d'ailleurs} $\|v\|_2 = \|W\|_2$ ($= \sqrt{\|u\|_1}$), \uncertain{et} \uncertain{aussi} \struck{$\|u\|_1$} $u = vW$. \ill{} \ill{} ; \ill{} $\|v\|_2\, \|W\|_2 \leq \|u\|_1$.
\note{la page s'arrête sur cette ligne}

\section{Produits extérieurs hilbertiens. Applications : des inégalités diverses}

\note{Titre de sa main, en tête de la page 58, numéroté « 3. ». La numérotation des formules repart de (1).}

\page{58}
\uncertain{Soient} $E_1, \ldots, E_n$ \uncertain{n} \uncertain{espaces} \uncertain{de} Hilbert. \uncertain{Alors} \uncertain{par} \uncertain{récurrence}, \uncertain{on} \ill{} \uncertain{leur} \uncertain{produit} \uncertain{tensoriel} \uncertain{hilbertien} $\bigotimes^{(2)}_{1 \leq i \leq n} E_i$, \uncertain{évidemment} \uncertain{l'espace} \uncertain{préhilbertien} \uncertain{séparé} \ill{}, \uncertain{à} \uncertain{une} \uncertain{structure} \uncertain{préhilbertienne} \uncertain{sur} \uncertain{le} \uncertain{tensoriel} \ill{}, \uncertain{définie} \uncertain{par} \uncertain{la}
\[
  \text{(1)}\qquad (a_1 \otimes \cdots \otimes a_n,\; b_1 \otimes \cdots \otimes b_n) = (a_1, b_1) \cdots (a_n, b_n)
\]
\uncertain{qui} \uncertain{résulte} \uncertain{aussitôt} \uncertain{du} \uncertain{N°} 1 \uncertain{qu'on} \uncertain{obtient} \uncertain{là} \uncertain{une} \uncertain{forme} \uncertain{sesquilinéaire} \ill{} \uncertain{positive}, \struck{\ill{}} \uncertain{bien} \uncertain{hilbertienne} \uncertain{et} \uncertain{donne} \ill{} \struck{\uncertain{c'est} \uncertain{l'espace} \uncertain{produit} \uncertain{tensoriel} \uncertain{hilbertien} \uncertain{des} $E_i$, \uncertain{et}} \note{ligne biffée et enfermée dans une boucle avec la précédente ; en interligne : « $\bigotimes^{(2)}_{i \in I_\alpha}$ \ill{} »} \uncertain{hilbertien} \ill{} \uncertain{de} \ill{} \ill{} : $\bigotimes^{(2)}_{i} E_i$ \uncertain{si} \ill{} \uncertain{les} $E_i$ \uncertain{sont} \uncertain{isomorphes} \ill{} \uncertain{à} $E$. \ill{} \ill{} \uncertain{Si} $(e^{i}_{\alpha_i})_{\alpha_i \in A_i}$ \uncertain{est} \uncertain{une} \uncertain{base} \uncertain{orthonormale} \uncertain{de} $E_i$ \uncertain{pour} \uncertain{tout} $i$, \uncertain{alors} \uncertain{on} \uncertain{obtient} \uncertain{une} \uncertain{base} \uncertain{orthonormale} $(e_\alpha)_{\alpha \in A}$ \uncertain{de} $\bigotimes^{(2)}_{1 \leq i \leq n} E_i$, \uncertain{où} $A = A_1 \times \cdots \times A_n$, \uncertain{et} \uncertain{pour} $\alpha = (\alpha_1, \ldots, \alpha_n)$, \uncertain{on} \uncertain{pose} $e_\alpha = e^{1}_{\alpha_1} \otimes \cdots \otimes e^{n}_{\alpha_n}$.

\uncertain{Soient} \uncertain{maintenant} \uncertain{deux} \uncertain{séquences} $(E_i)$, $(F_i)_{1 \leq i \leq n}$ \uncertain{d'espaces} \uncertain{de} Hilbert, \uncertain{et} \uncertain{soit} \uncertain{pour} \uncertain{tout} $i$, $u_i$ \uncertain{une} \uncertain{application} \uncertain{linéaire} \uncertain{continue} \uncertain{de} $E_i$ \uncertain{dans} $F_i$, \uncertain{alors} \uncertain{l'application} $\bigotimes u_i$ \uncertain{de} $\bigotimes_{1 \leq i \leq n} E_i$ \uncertain{dans} $\bigotimes F_i$ \uncertain{est} \uncertain{continue} \uncertain{pour} \uncertain{les} \uncertain{structures} \uncertain{hilbertiennes}, \uncertain{et} \emph{\uncertain{de} \uncertain{façon} \uncertain{précise}}
\[
  \text{(2)}\qquad \|u_1 \otimes \cdots \otimes u_n\| \leq \|u_1\| \cdots \|u_n\|
\]
\uncertain{Il} \uncertain{suffit} \uncertain{de} \uncertain{le} \uncertain{démontrer} \uncertain{quand} \uncertain{tous} \uncertain{les} $u_i$ \uncertain{sauf} \uncertain{un} \uncertain{sont} \uncertain{l'identité}, \struck{\uncertain{et} \uncertain{alors}} \uncertain{le} \uncertain{résultat} \uncertain{en} \uncertain{découle} \ill{}
\note{la phrase se poursuit page 59}

\marginal{en marge gauche, sideways, sur plusieurs lignes : « \emph{l'associativité de l'opération} \ill{} ; $1 \leq \alpha < \beta \cdots < \gamma < n$, \ill{} ; \uncertain{avec} $(\bigotimes^{(2)}_{1 \leq i \leq \alpha} E_i) \otimes^{(2)} (\bigotimes^{(2)}_{\alpha < i \leq \beta} E_i) \otimes \cdots \otimes (\bigotimes^{(2)}_{\gamma < i \leq n} E_i)$ \ill{} ; \ill{} \uncertain{est} \uncertain{aussi} \uncertain{le} \uncertain{produit} \uncertain{tensoriel} \uncertain{hilbertien} : \uncertain{si} \ill{} \uncertain{de} $\bigotimes^{(2)} E_i$ \uncertain{à} \ill{} »}

\page{59}
\struck{\uncertain{de} \uncertain{le} \uncertain{prouver} \uncertain{au} \uncertain{produit}, \uncertain{en} \uncertain{suivant} \ill{}} \uncertain{car} \uncertain{on} \uncertain{a} $u_1 \otimes \cdots \otimes u_n = U_1 \circ \cdots \circ U_n$, \uncertain{où} $U_i$ \uncertain{est} \uncertain{le} \struck{$u = 1 \otimes \cdots \otimes 1 \otimes \ill{}$} \note{ligne lourdement biffée} \uncertain{produit} \uncertain{tensoriel} \uncertain{de} $n$ \uncertain{opérateurs}, \uncertain{tous} \uncertain{égaux} \uncertain{à} \uncertain{l'identité} \uncertain{sauf} \uncertain{le} $i$-\uncertain{ème} \uncertain{égal} \uncertain{à} $u_i$. \struck{\uncertain{Mais} \ill{} \uncertain{Prenons} \uncertain{par} \ill{}} \uncertain{Mais}, \struck{$u \otimes 1 \otimes \cdots \otimes 1$}, \uncertain{où} $A \in L(E_1, F_1)$, \uncertain{la} \uncertain{première} \uncertain{façon} \uncertain{la} \uncertain{plus} \uncertain{simple} \ill{} \uncertain{identifions} $E_i$ \ill{} \uncertain{avec} $l^{2}(A_i)$, \uncertain{alors} $F_i$ \ill{} \uncertain{avec} \ill{}, $\bigotimes^{(2)}_{1 \leq i \leq n} E_i$ \uncertain{est} \uncertain{identifié} \uncertain{à} \ill{} $l^{2}(A_1 \times \cdots \times A_n)$ \ill{} $E_1$ \uncertain{et} \uncertain{identifié} \ill{} $l^{2}(A_1 \ldots$ \ill{} $F_1 \otimes^{(2)} E_2 \otimes^{(2)} \cdots \otimes^{(2)} E_n$ \note{les cinq lignes de « $A \in L(E_1,F_1)$ » à « $F_1 \otimes E_2 \otimes \cdots$ » sont enfermées dans une longue boucle} \ill{} \uncertain{en} \uncertain{vertu} \uncertain{de} \uncertain{l'associativité} \uncertain{du} \uncertain{produit} \uncertain{tensoriel} \uncertain{hilbertien}, \uncertain{on} \uncertain{est} \uncertain{ramené} \uncertain{au} \uncertain{cas} \uncertain{de} \uncertain{deux} \uncertain{facteurs} $E_1$ \uncertain{et} $E_2$, $F_1$ \uncertain{et} $F_2 = E_2$, \uncertain{et} \uncertain{à} $A_1 \in L(E_1, F_1)$, $A_2 = 1$. \uncertain{Mais} \uncertain{alors}, \uncertain{identifiant} $E_1 \otimes^{(2)} E_2$ \uncertain{à} $L^{(2)}(\bar{E}_1, E_2)$, $F_1 \otimes^{(2)} E_2$ \uncertain{à} $L^{(2)}(\bar{F}_1, E_2)$, \uncertain{on} \uncertain{est} \uncertain{ramené} \uncertain{à} \uncertain{prouver} \uncertain{que} \uncertain{pour} \uncertain{tout} $u \in L^{(2)}(\bar{E}_1, E_2)$, \uncertain{on} \uncertain{a} $u \circ {}^{t}\!A \in L^{(2)}(\bar{F}_1, E_2)$ \note{lecture incertaine de l'opérateur composé avec $u$ ; un $A$ affecté d'un signe, transposé ou adjoint}, \uncertain{et} $\|u\, {}^{t}\!A\|_2 \leq \|u\|_2\, \|A\|$, \uncertain{ce} \uncertain{qui} \uncertain{n'est} \uncertain{en} \uncertain{effet} \uncertain{que} \uncertain{la} \uncertain{formule} (17) \uncertain{du} N° 1 \note{ainsi sur la page ; la formule (17) est celle du N° 2 ci-dessus}. \uncertain{Par} \uncertain{suite}, \struck{\uncertain{l'application} $u_1 \otimes \cdots \otimes$} \uncertain{Par} \uncertain{suite} $\bigotimes u_i$ \struck{\ill{}} \uncertain{se} \struck{\ill{}} \uncertain{prolonge} \uncertain{en} \uncertain{une} \uncertain{application} \uncertain{de} $\bigotimes^{(2)}_{1 \leq i \leq n} E_i$ \uncertain{dans} $\bigotimes^{(2)} F_i$, \uncertain{de} \uncertain{norme} $\leq$ \ill{} (\uncertain{indiquée} \uncertain{dans} \ill{}), \uncertain{de} \uncertain{la} \uncertain{même} \uncertain{façon}. \uncertain{En} \emph{\uncertain{particulier}} \uncertain{ceci} \uncertain{s'applique} \uncertain{lorsque} \uncertain{tous} \uncertain{les} $E_i$ \uncertain{sont} \uncertain{un} \uncertain{même} \uncertain{espace} \uncertain{de} Hilbert $E$, \uncertain{considérons} \struck{\ill{} \uncertain{espaces} \uncertain{de} \uncertain{Hilbert} $E_i \ldots E_n$} $\bigotimes^{(2)} E$ \struck{\uncertain{Posons} \uncertain{encore} \ill{}} \uncertain{Posons} \ill{}
\[
  \text{(3)}\qquad a_n = \frac{1}{n!} \sum_{\sigma \in \mathfrak{S}_n} \varepsilon_\sigma\, \sigma
\]
\note{le $\varepsilon_\sigma$ est marqué au crayon rouge ; le groupe symétrique est un S gothique} \uncertain{où} \uncertain{le} \uncertain{second} \uncertain{membre} \uncertain{est} \uncertain{considéré} \uncertain{comme} \uncertain{opérateur} \uncertain{dans} $\bigotimes^{n\,(2)} E$ \note{« $\otimes$ » avec « $n$ » et « (2) » en exposant} ($\varepsilon_\sigma$ \uncertain{signe} \uncertain{de} \uncertain{la} \uncertain{permutation} $\sigma$). \struck{\ill{}} \uncertain{au} \uncertain{calcul}, \uncertain{on} \uncertain{trouve} \ill{} :
\note{la phrase se poursuit page 60}

\page{60}
\ill{}, \uncertain{l'opérateur} $a_n$ \note{« $a_n$ » en interligne} \uncertain{est} \ill{} \uncertain{dans} $\bigotimes^{(2)}_{1 \leq i \leq n} E_i$ \struck{\ill{}} \struck{\ill{} \uncertain{que} $a_n$ \uncertain{est}} \uncertain{évidemment} \uncertain{un} \emph{\uncertain{projecteur}} \uncertain{hermitien} (\uncertain{car} \uncertain{il} \uncertain{est} \uncertain{hermitien}, \uncertain{et} \uncertain{on} \uncertain{sait} \uncertain{que} \struck{\ill{} $a_n^{2} = a_n$} \uncertain{on} \uncertain{a} \uncertain{la} \struck{\ill{}} \uncertain{propriété} \uncertain{connue} \uncertain{algébrique} \ill{} \uncertain{formule}, \uncertain{par} \uncertain{définition}, \uncertain{qui} \uncertain{est} \uncertain{dense}), \uncertain{d'ailleurs} \uncertain{hermitien} (\uncertain{noter} \uncertain{que} \uncertain{l'adjoint} \uncertain{de} \uncertain{l'opérateur} \uncertain{de} $\sigma$ \uncertain{est} $\sigma^{-1}$, \uncertain{à} \uncertain{cause} \uncertain{de} \uncertain{la} \uncertain{formule}
\[
  (a_{\sigma 1} \otimes \cdots \otimes a_{\sigma n},\; b_1 \otimes \cdots \otimes b_n) = (a_1 \otimes \cdots \otimes a_n,\; b_{\sigma^{-1} 1} \otimes \cdots \otimes b_{\sigma^{-1} n})
\]
\uncertain{conséquence} \uncertain{immédiate} \uncertain{de} (1)). \struck{\ill{}}

\emph{Définition.} \uncertain{Soit} $E$ \uncertain{un} \uncertain{espace} \uncertain{de} Hilbert, $n$ \uncertain{un} \uncertain{entier} $\geq 0$, \uncertain{on} \uncertain{appelle} \uncertain{puissance} \uncertain{extérieure} \uncertain{hilbertienne} $n$-\uncertain{ième} \uncertain{de} $E$ \note{en interligne : « l'espace de Hilbert »}, \uncertain{et} \uncertain{on} \uncertain{note} $\Lambda^{n\,(2)} E$ \note{$\Lambda$ avec « $n$ » et « (2) » en exposant ; accolade en marge gauche sur la définition}, \struck{\ill{}} \uncertain{le} \uncertain{sous-espace} \uncertain{de} $\bigotimes^{n\,(2)} E$ \uncertain{formé} \uncertain{des} \ill{} \uncertain{l'opérateur} $a_n$ \uncertain{défini} \uncertain{par} (3), \uncertain{muni} \uncertain{du} \uncertain{produit} \uncertain{scalaire} \uncertain{et} \uncertain{multiplié} \ill{} \uncertain{par} $n!$ \ill{} \uncertain{Comme} \uncertain{on} \uncertain{a} \uncertain{aussi} \uncertain{la} \ill{} \uncertain{sous-espace} (\uncertain{ce} \uncertain{facteur} $n!$ \uncertain{près}) \uncertain{de} $\bigotimes^{(2)} E$ \uncertain{image} \uncertain{du} \uncertain{projecteur} $a_n$. \uncertain{Pour} \uncertain{définition} \ill{} \uncertain{si} \uncertain{on} \uncertain{veut} \uncertain{qu'il} \uncertain{y} \uncertain{a} \uncertain{une} \uncertain{relation} \emph{\uncertain{biunivoque}} \uncertain{pour} \ill{} \uncertain{naturelle} \uncertain{la} \uncertain{puissance} \uncertain{extérieure} \uncertain{algébrique} $\Lambda E$ \uncertain{dans} $\Lambda^{(2)} E$, \uncertain{dont} \uncertain{l'image} \uncertain{est} \uncertain{identique} \uncertain{à} \uncertain{l'image} \uncertain{de} $\bigotimes E$ \uncertain{dans} $\bigotimes^{(2)} E$. \uncertain{On} \uncertain{identifie} \uncertain{donc} $\Lambda E$ \uncertain{à} \uncertain{une} \uncertain{partie} \uncertain{de} $\Lambda^{(2)} E$, \uncertain{dense} \uncertain{dans} \uncertain{ce} \uncertain{dernier} \ill{} \uncertain{de} \uncertain{plus} \emph{\uncertain{dense}} \uncertain{dans} $\Lambda^{(2)} E$. \ill{} \ill{} \uncertain{préhilbertien} \ill{} \uncertain{est} \uncertain{le} \uncertain{complété} \uncertain{de} $\Lambda E$ \ill{} $\Lambda E$ \uncertain{s'explicite} \uncertain{aisément} \uncertain{sur} \uncertain{les} \uncertain{multivecteurs} \uncertain{décomposables} \note{trait ondulé au crayon rouge sous cette ligne} : \uncertain{Posons} $A = a_1 \otimes \cdots \otimes a_n$, $B = b_1 \otimes \cdots \otimes b_n$, \uncertain{on} \uncertain{a}
\[
  \struck{(a_1 \wedge \cdots \wedge a_n,\; b_1 \wedge \cdots \wedge b_n) = \Bigl(\frac{1}{n!}\Bigr)^{2} \sum_{\sigma \in \mathfrak{S}_n} \varepsilon_\sigma\, \sigma\, a_1 \ldots}
\]
\note{ligne biffée d'un trait, le facteur $(1/n!)^{2}$ cerné}
\[
  (a_1 \wedge \cdots \wedge a_n,\; b_1 \wedge \cdots \wedge b_n) = n!\,(a_n A,\; a_n B) = n!\,(A,\; a_n^{2} B) = n!\,(A,\; a_n B)
\]
\[
  = \sum_{\sigma \in \mathfrak{S}_n} \varepsilon_\sigma\, (a_1, b_{\sigma 1}) \cdots (a_n, b_{\sigma n}) = \det\bigl((a_i, b_j)\bigr), \qquad \text{donc}
\]
\ill{} \struck{\ill{}} \ill{} (
\note{la page, et le lot, s'arrêtent sur cette parenthèse ouverte}

\end{document}
